% =============================================================================
% L01 (Class 1, Review of basic ingredients) --- book exercises, verbatim statements + solutions.
% Built from the MPE exam template (lista-exam.tex). One source, two outputs:
%   \solucoesfalse  -> blank list (space to work)
%   \solucoestrue   -> answer key (same file)
% Compile: pdflatex (2x); check: pdffonts <file>.pdf (Type 1 only).
% Every number printed in the solutions is reproduced and asserted by
%   python Listas/L01-book-exercises-check.py
% =============================================================================
\documentclass[11pt,a4paper]{article}
\usepackage[utf8]{inputenc}
\usepackage[T1]{fontenc}
\usepackage{lmodern}   % Type 1 vector fonts (replaces PK bitmaps)
\usepackage{cmap}      % ToUnicode CMap -> copyable text and accents
\usepackage[brazil,english]{babel}
\usepackage{amsmath,amssymb,amsfonts}
\usepackage[a4paper,margin=2.2cm]{geometry}
\usepackage{enumitem}
\usepackage{fancyhdr}
\usepackage{comment}
% One cut file per job: parallel builds in one folder would otherwise share comment.cut and
% swap solutions. Redefine DefaultCutFileName; overriding CommentCutFile alone is reset by comment.sty.
\def\DefaultCutFileName{\def\CommentCutFile{\jobname.cut}}
\usepackage{needspace}
\usepackage{xcolor}
\usepackage{microtype}
\usepackage{booktabs}
\usepackage{pgfplots}
\pgfplotsset{compat=1.16}

% ====================== CONFIGURATION ======================
\newif\ifsolucoes
\solucoesfalse                 % <<<< \solucoestrue for the ANSWER KEY; \solucoesfalse for the LIST

\newcommand{\disciplina}{Econometria Avançada --- MPE (Insper)}
\newcommand{\listatema}{\textit{L01 --- Review of basic ingredients (book exercises)}}
\newcommand{\espacoresol}{4.5cm}   % height of the working space in the blank list

% ====================== MACHINERY ======================
\ifsolucoes\includecomment{sol}\else\excludecomment{sol}\fi

\pagestyle{fancy}\fancyhf{}
\lhead{\footnotesize\disciplina}
\rhead{\footnotesize L01 book exercises\ifsolucoes{} --- Solutions\fi}
\cfoot{\footnotesize\thepage}
\renewcommand{\headrulewidth}{0.3pt}
\setlength{\headheight}{14pt}

\setlist{topsep=2pt,itemsep=2pt}
\newcounter{questao}
\numberwithin{equation}{questao}   % equation (q.m): question q, equation m

% Block header (one block per book, in the class plan's order)
\newcommand{\bloco}[3]{\bigskip\noindent
  {\large\bfseries\color{blue!55!black}Block #1 --- #2}\hfill\textit{(#3)}\par
  \noindent\rule{\linewidth}{0.3pt}\smallskip}

% Question (automatic numbering)
\newcommand{\questao}[1]{\refstepcounter{questao}\medskip\needspace{5\baselineskip}\noindent
  \textbf{Question \thequestao.}~#1\par}
% Source line that opens every book statement
\newcommand{\fonte}[1]{\par\noindent{\small\textbf{Source:} #1}\par\smallskip}
% Notation list that opens every solution (estilo-avaliacao §4, "Equations, notation and layout").
\newenvironment{notacao}{\par\smallskip\noindent\textbf{Notation.}\begin{itemize}[leftmargin=1.4em,topsep=1pt,itemsep=0pt]\small}{\end{itemize}}
\allowdisplaybreaks

% Working space in the blank list
\newcommand{\resolespaco}{\ifsolucoes\else\par\textit{Solution:}\vspace{\espacoresol}\par\fi}

% Solution step (use inside sol)
\newcommand{\passo}[2]{\par\smallskip\noindent\textbf{Step --- #1.}~#2}
\newcommand{\intuicao}[1]{\par\smallskip\noindent\textit{Intuition.}~#1}
\newcommand{\refbib}[1]{\par\smallskip\noindent{\footnotesize\color{black!60}\textbf{Ref:} #1}}
\newcommand{\solhead}{\par\medskip\noindent{\color{blue!55!black}\rule{\linewidth}{0.3pt}}\par\noindent\textbf{Solution.}}

\newcommand{\E}{\mathbb{E}}
\newcommand{\Var}{\operatorname{Var}}
\newcommand{\Cov}{\operatorname{Cov}}

% =============================================================================
\begin{document}
\begin{center}
  {\Large\bfseries \disciplina}\\[2pt]
  {\large Exercise list: \listatema}\ifsolucoes\\[2pt]\textit{Solutions}\fi
\end{center}
\vspace{2pt}

{\footnotesize\noindent\textbf{Instructions.} Every question is a textbook exercise from the
\emph{Practice} list of Class~1 in \texttt{Map/class-plan.md}, reproduced verbatim from the book PDF
(statements from Portuguese-language books stay in Portuguese). Show every step. Computer exercises
are solved in Python; the script \texttt{Listas/L01-book-exercises-check.py} reproduces and asserts
every number printed in the solutions. Book abbreviations follow the class plan: T = Tsay, W =
Wooldridge, G = Greenberg, I = ISLP, B = Bolfarine \& Sandoval, MS = Morettin \& Singer.\par}

% =============================================================================
\bloco{1}{Tsay, \textit{Analysis of Financial Time Series}}{T 2.2}

\questao{T 2.2 --- AR(2) moments and forecasts.}
\fonte{Tsay, R.~S. (2010), \textit{Analysis of Financial Time Series}, 3rd ed., Wiley, Chapter~2
Exercises, Exercise~2.2, p.~104 --- reproduced verbatim. (Header of the exercise set: ``If not
specifically specified, use 5\% significance level to draw conclusions in the exercises.'')}

2.2. Suppose that the daily log return of a security follows the model
\[ r_t = 0.01 + 0.2 r_{t-2} + a_t , \]
where $\{a_t\}$ is a Gaussian white noise series with mean zero and variance 0.02. What are the mean
and variance of the return series $r_t$? Compute the lag-1 and lag-2 autocorrelations of $r_t$.
Assume that $r_{100} = -0.01$, and $r_{99} = 0.02$. Compute the 1- and 2-step-ahead forecasts of the
return series at the forecast origin $t = 100$. What are the associated standard deviations of the
forecast errors?
\resolespaco
\begin{sol}
\solhead
\begin{notacao}
\item $r_t$: daily log return; $a_t\sim N(0,\sigma_a^2)$ i.i.d.\ shock, $\sigma_a^2=0.02$.
\item $\phi_0=0.01$, $\phi_2=0.2$ (the lag-1 coefficient is $\phi_1=0$).
\item $\mu=\E(r_t)$; $\gamma_\ell=\Cov(r_t,r_{t-\ell})$; $\rho_\ell=\gamma_\ell/\gamma_0$, $\ell\ge 0$.
\item $\hat r_{100}(\ell)=\E(r_{100+\ell}\mid r_{100},r_{99},\dots)$: $\ell$-step forecast at origin 100;
  $e_{100}(\ell)=r_{100+\ell}-\hat r_{100}(\ell)$: forecast error.
\end{notacao}

\passo{Model and stationarity}{The model is}
\begin{equation}\label{q1:model}
r_t=\phi_0+\phi_2 r_{t-2}+a_t .
\end{equation}
Its characteristic equation $1-\phi_2x^2=0$ has roots $x=\pm 1/\sqrt{\phi_2}=\pm 1/\sqrt{0.2}=\pm 2.2361$,
both outside the unit circle, so $r_t$ is weakly stationary and $\mu$, $\gamma_\ell$ do not depend on $t$.

\passo{Mean}{Taking expectations of \eqref{q1:model} and using $\E(a_t)=0$:}
\begin{equation}\label{q1:meaneq}
\mu=\phi_0+\phi_2\mu .
\end{equation}
Solving \eqref{q1:meaneq} for $\mu$:
\begin{equation}\label{q1:mean}
\mu=\frac{\phi_0}{1-\phi_2}.
\end{equation}
Substituting the numbers into \eqref{q1:mean}:
\[
\mu=\frac{(0.01)}{1-(0.2)}=\frac{0.01}{0.8}=0.0125 .
\]

\passo{Autocovariance equations}{Subtracting \eqref{q1:meaneq} from \eqref{q1:model}:}
\begin{equation}\label{q1:dev}
r_t-\mu=\phi_2\,(r_{t-2}-\mu)+a_t .
\end{equation}
Multiplying \eqref{q1:dev} by $(r_{t-\ell}-\mu)$ and taking expectations:
\begin{equation}\label{q1:yw}
\gamma_\ell=\phi_2\,\gamma_{\ell-2}+\E\big[a_t(r_{t-\ell}-\mu)\big],
\qquad
\E\big[a_t(r_{t-\ell}-\mu)\big]=\begin{cases}\sigma_a^2,&\ell=0,\\ 0,&\ell\ge1,\end{cases}
\end{equation}
because $a_t$ is uncorrelated with the past and $\E[a_t(r_t-\mu)]=\E(a_t^2)=\sigma_a^2$ by \eqref{q1:dev}.

\passo{Variance}{Setting $\ell=0$ and $\ell=2$ in \eqref{q1:yw}, with $\gamma_{-2}=\gamma_2$:}
\begin{align}
\gamma_0&=\phi_2\gamma_2+\sigma_a^2, \label{q1:g0}\\
\gamma_2&=\phi_2\gamma_0. \label{q1:g2}
\end{align}
Substituting \eqref{q1:g2} into \eqref{q1:g0}:
\[
\gamma_0=\phi_2\,(\phi_2\gamma_0)+\sigma_a^2 .
\]
Collecting $\gamma_0$:
\begin{equation}\label{q1:var}
\gamma_0=\frac{\sigma_a^2}{1-\phi_2^2}.
\end{equation}
Substituting the numbers into \eqref{q1:var}:
\[
\gamma_0=\frac{(0.02)}{1-(0.2)^2}=\frac{0.02}{0.96}=0.020833,
\qquad \sqrt{\gamma_0}=0.1443 .
\]

\passo{Autocorrelations}{Setting $\ell=1$ in \eqref{q1:yw}, with $\gamma_{-1}=\gamma_1$:}
\begin{equation}\label{q1:g1}
\gamma_1=\phi_2\gamma_1 \;\Longrightarrow\; (1-\phi_2)\gamma_1=0 \;\Longrightarrow\; \gamma_1=0,
\end{equation}
since $\phi_2=0.2\neq1$. Hence $\rho_1=\gamma_1/\gamma_0=0$. Dividing \eqref{q1:g2} by $\gamma_0$:
\[
\rho_2=\frac{\gamma_2}{\gamma_0}=\phi_2=0.2 .
\]

\passo{1-step forecast}{Writing \eqref{q1:model} at $t=101$:}
\begin{equation}\label{q1:r101}
r_{101}=\phi_0+\phi_2 r_{99}+a_{101}.
\end{equation}
$r_{99}$ is known at the origin and $\E(a_{101}\mid\cdot)=0$, so
\begin{equation}\label{q1:f1}
\hat r_{100}(1)=\phi_0+\phi_2 r_{99}.
\end{equation}
Substituting the numbers into \eqref{q1:f1}:
\[
\hat r_{100}(1)=0.01+0.2\,(0.02)=0.01+0.004=0.014 .
\]

\passo{2-step forecast}{Writing \eqref{q1:model} at $t=102$:}
\begin{equation}\label{q1:r102}
r_{102}=\phi_0+\phi_2 r_{100}+a_{102},
\end{equation}
and $r_{100}$ is known, so
\begin{equation}\label{q1:f2}
\hat r_{100}(2)=\phi_0+\phi_2 r_{100}.
\end{equation}
Substituting the numbers into \eqref{q1:f2}:
\[
\hat r_{100}(2)=0.01+0.2\,(-0.01)=0.01-0.002=0.008 .
\]

\passo{Forecast-error standard deviations}{Subtracting \eqref{q1:f1} from \eqref{q1:r101} and
\eqref{q1:f2} from \eqref{q1:r102}:}
\begin{equation}\label{q1:err}
e_{100}(1)=a_{101},\qquad e_{100}(2)=a_{102}.
\end{equation}
Both have variance $\sigma_a^2$, so
\[
\operatorname{sd}[e_{100}(1)]=\operatorname{sd}[e_{100}(2)]=\sqrt{0.02}=0.1414 .
\]
(In a general AR(2), $e_h(2)=a_{h+2}+\phi_1a_{h+1}$; here $\phi_1=0$, so the 2-step error is as small as
the 1-step error.)

\intuicao{With $\phi_1=0$ the series is two interleaved AR(1) processes, one on odd and one on even
days. Returns two days apart are correlated ($\rho_2=0.2$), adjacent days are not ($\rho_1=0$). The
2-step forecast is driven by a return already observed ($r_{100}$), so it is as precise as the
1-step one; both forecasts sit close to $\mu=0.0125$ because $\phi_2$ is small, and the error
standard deviation $0.1414$ is only slightly below the unconditional $0.1443$.}
\refbib{Tsay (2010), §2.4.1 (properties of AR models, pp.~38--46) and §2.4.4 (forecasting, pp.~54--56).}
\end{sol}

% =============================================================================
\bloco{2}{Wooldridge, \textit{Introductory Econometrics}}{W 11.C8; W 12.1}

\questao{W 11.C8 --- AR(1) for unemployment (computer exercise).}
\fonte{Wooldridge, J.~M. (2016), \textit{Introductory Econometrics: A Modern Approach}, 6th ed.,
Cengage, Chapter~11, Computer Exercise C8, p.~370 --- reproduced verbatim.}

C8 Use the data in PHILLIPS for this exercise.
\begin{enumerate}[label=(\roman*),leftmargin=2.4em]
\item Estimate an AR(1) model for the unemployment rate. Use this equation to predict the
unemployment rate for 2004. Compare this with the actual unemployment rate for 2004. (You can find
this information in a recent \textit{Economic Report of the President}.)
\item Add a lag of inflation to the AR(1) model from part (i). Is $inf_{t-1}$ statistically
significant?
\item Use the equation from part (ii) to predict the unemployment rate for 2004. Is the result better
or worse than in the model from part (i)?
\item Use the method from Section 6-4 to construct a 95\% prediction interval for the 2004
unemployment rate. Is the 2004 unemployment rate in the interval?
\end{enumerate}
\resolespaco
\begin{sol}
\solhead
\begin{notacao}
\item $unem_t$: U.S.\ civilian unemployment rate in year $t$ (percent); $inf_t$: CPI inflation rate
  in year $t$ (percent). PHILLIPS covers 1948--2003; with one lag the estimation sample is
  1949--2003, $n=55$.
\item $\beta_j$: population coefficients; $\hat\beta_j$: OLS estimates; $\operatorname{se}(\cdot)$:
  usual OLS standard error; $u_t$: error; $\hat\sigma^2=\mathrm{SSR}/(n-k-1)$; $R^2$.
\item $\hat y_0$: prediction of $unem_{2004}$; $\hat e_0=unem_{2004}-\hat y_0$: prediction error;
  $\theta_0$: population value of the prediction. Actual 2004 rate: $5.5$ (BLS annual average,
  as reported in the \textit{Economic Report of the President}; the mean of the 12 monthly FRED
  \texttt{UNRATE} values for 2004 is $5.54$).
\item 2003 values used for the 2004 prediction: $unem_{2003}=6.0$, $inf_{2003}=2.3$.
\end{notacao}

\passo{(i) AR(1) estimates}{The model is}
\begin{equation}\label{q2:ar1}
unem_t=\beta_0+\beta_1\,unem_{t-1}+u_t .
\end{equation}
OLS on 1949--2003 gives
\begin{equation}\label{q2:ar1hat}
\widehat{unem}_t=\underset{(0.5202)}{1.489685}+\underset{(0.0893)}{0.742382}\,unem_{t-1},
\qquad n=55,\quad R^2=0.5660,\quad \hat\sigma=0.9986 .
\end{equation}

\passo{(i) Prediction for 2004}{Evaluating \eqref{q2:ar1hat} at $unem_{2003}$:}
\begin{equation}\label{q2:f1}
\widehat{unem}_{2004}=\hat\beta_0+\hat\beta_1\,unem_{2003}.
\end{equation}
Substituting the numbers into \eqref{q2:f1}:
\[
\widehat{unem}_{2004}=1.489685+0.742382\,(6.0)=1.489685+4.454292=5.9440 .
\]
Prediction error against the actual rate:
\[
\hat y_0-unem_{2004}=5.9440-5.5=0.4440 ,
\]
so the AR(1) overpredicts 2004 unemployment by about 0.44 percentage points.

\passo{(ii) Adding lagged inflation}{The model is}
\begin{equation}\label{q2:ar1inf}
unem_t=\beta_0+\beta_1\,unem_{t-1}+\beta_2\,inf_{t-1}+u_t ,
\end{equation}
with OLS estimates
\begin{equation}\label{q2:ar1infhat}
\widehat{unem}_t=\underset{(0.4413)}{1.296416}+\underset{(0.0779)}{0.649479}\,unem_{t-1}
+\underset{(0.0388)}{0.183009}\,inf_{t-1},
\qquad n=55,\quad R^2=0.6963 .
\end{equation}
The $t$ statistic for $H_0:\beta_2=0$ is
\begin{equation}\label{q2:t}
t_{\hat\beta_2}=\frac{\hat\beta_2}{\operatorname{se}(\hat\beta_2)}=\frac{0.183009}{0.0388}=4.72 ,
\end{equation}
against a two-sided 5\% critical value $t_{52,\,0.975}=2.0066$ (two-sided $p$-value $0.00002$).
So $inf_{t-1}$ is statistically significant at any conventional level. The positive sign says that
a year of higher inflation is followed by higher unemployment, the pattern of an
expectations-augmented Phillips curve with supply shocks, not a stable inflation--unemployment
trade-off.

\passo{(iii) Prediction from \eqref{q2:ar1infhat}}{Evaluating \eqref{q2:ar1infhat} at the 2003 values:}
\begin{equation}\label{q2:f2}
\widehat{unem}_{2004}=\hat\beta_0+\hat\beta_1\,unem_{2003}+\hat\beta_2\,inf_{2003}.
\end{equation}
Substituting the numbers into \eqref{q2:f2}:
\[
\widehat{unem}_{2004}=1.296416+0.649479\,(6.0)+0.183009\,(2.3)
=1.296416+3.896874+0.420921=5.6142 .
\]
Prediction error:
\[
\hat y_0-unem_{2004}=5.6142-5.5=0.1142 .
\]
Since $0.1142<0.4440$, the model with lagged inflation predicts 2004 better (for this one year).

\passo{(iv) Section 6-4: the prediction and its standard error}{The population prediction is}
\begin{equation}\label{q2:theta0}
\theta_0=\beta_0+\beta_1(6.0)+\beta_2(2.3).
\end{equation}
Solving \eqref{q2:theta0} for $\beta_0$ and substituting into \eqref{q2:ar1inf}:
\begin{equation}\label{q2:centred}
unem_t=\theta_0+\beta_1\,(unem_{t-1}-6.0)+\beta_2\,(inf_{t-1}-2.3)+u_t ,
\end{equation}
so the OLS intercept of \eqref{q2:centred} is $\hat y_0=5.6142$ (the same number as in (iii)) and its
standard error is $\operatorname{se}(\hat y_0)=0.136381$.

\passo{(iv) Standard error of the prediction error}{The prediction error
$\hat e_0=u_{2004}-(\hat y_0-\theta_0)$ adds the new error to the estimation error, which are
uncorrelated, so (Wooldridge eq.~(6.36))}
\begin{equation}\label{q2:see}
\operatorname{se}(\hat e_0)=\sqrt{\operatorname{se}(\hat y_0)^2+\hat\sigma^2}.
\end{equation}
With $\hat\sigma^2=\mathrm{SSR}/52=0.711321$, substituting into \eqref{q2:see}:
\[
\operatorname{se}(\hat e_0)=\sqrt{(0.136381)^2+(0.711321)}=\sqrt{0.018600+0.711321}
=\sqrt{0.729921}=0.8544 .
\]

\passo{(iv) 95\% prediction interval}{The interval (Wooldridge eq.~(6.37)) is}
\begin{equation}\label{q2:pi}
\hat y_0\pm t_{52,\,0.975}\cdot\operatorname{se}(\hat e_0).
\end{equation}
Substituting into \eqref{q2:pi}: the half-width is $2.0066\,(0.8544)=1.7144$, so
\[
[\,5.6142-1.7144,\;5.6142+1.7144\,]=[\,3.8998,\;7.3286\,].
\]
The actual 2004 rate, $5.5$, lies inside the interval.

\passo{Python}{Code and key output (the full version, with assertions, is in the check script):}
{\small
\begin{verbatim}
import wooldridge, statsmodels.api as sm, numpy as np
ph = wooldridge.data("phillips").sort_values("year")
ph["unem_1"], ph["inf_1"] = ph.unem.shift(1), ph.inf.shift(1)
e = ph.dropna(subset=["unem_1", "inf_1"])                     # 1949-2003, n = 55
m1 = sm.OLS(e.unem, sm.add_constant(e[["unem_1"]])).fit()     # 1.489685, 0.742382
m2 = sm.OLS(e.unem, sm.add_constant(e[["unem_1", "inf_1"]])).fit()
#   const 1.296416 (0.4413)  unem_1 0.649479 (0.0779)  inf_1 0.183009 (0.0388), t = 4.72
X = np.column_stack([np.ones(55), e.unem_1 - 6.0, e.inf_1 - 2.3])
m3 = sm.OLS(e.unem.values, X).fit()                           # 5.6142, se 0.136381
se_e0 = np.sqrt(m3.bse[0]**2 + m3.scale)        # 0.8544 -> PI [3.8998, 7.3286]
\end{verbatim}
}

\intuicao{Unemployment is persistent ($\hat\beta_1\approx0.74$), so an AR(1) mostly extrapolates 2003's
6.0\%. Lagged inflation carries extra information about where the cycle is heading, which lowers
the 2004 forecast toward the realised 5.5\%. The prediction interval is wide because its width is
dominated by $\hat\sigma^2=0.711$, the variance of next year's shock; parameter uncertainty
($0.0186$) is a small part of it. One year of forecast comparison is weak evidence: a fair comparison
would use many out-of-sample years.}
\refbib{Wooldridge (2016), §11-2 and Example 11.5 (dynamic models, PHILLIPS), §18-5 (forecasting),
§6-4a (prediction intervals, pp.~186--190).}
\end{sol}

\questao{W 12.1 --- OLS standard errors under AR(1) errors.}
\fonte{Wooldridge, J.~M. (2016), \textit{Introductory Econometrics: A Modern Approach}, 6th ed.,
Cengage, Chapter~12, Problem~1, p.~396 --- reproduced verbatim.}

1 When the errors in a regression model have AR(1) serial correlation, why do the OLS standard
errors tend to underestimate the sampling variation in the $\hat\beta_j$? Is it always true that the
OLS standard errors are too small?
\resolespaco
\begin{sol}
\solhead
\begin{notacao}
\item $y_t=\beta_0+\beta_1x_t+u_t$, $t=1,\dots,n$; $x_t$ treated as nonrandom (condition on $X$),
  with $\bar x=0$ for simplicity (as in Wooldridge §12-1b).
\item $u_t=\rho u_{t-1}+e_t$, $|\rho|<1$, $e_t$ i.i.d.\ $(0,\sigma_e^2)$;
  $\sigma^2=\Var(u_t)=\sigma_e^2/(1-\rho^2)$.
\item $\mathrm{SST}_x=\sum_t x_t^2$; $\lambda$: first-order autocorrelation of $x_t$ when $x_t$ is
  itself AR(1) (used only for the illustration).
\end{notacao}

\passo{Sampling error of the slope}{With $\bar x=0$, OLS gives}
\begin{equation}\label{q3:b1}
\hat\beta_1=\beta_1+\frac{\sum_{t=1}^n x_tu_t}{\mathrm{SST}_x}.
\end{equation}
Under the AR(1) model the error covariances are
\begin{equation}\label{q3:cov}
\Cov(u_t,u_{t+j})=\rho^{j}\sigma^2,\qquad j\ge0 .
\end{equation}

\passo{True variance}{Taking the variance of \eqref{q3:b1}:}
\begin{equation}\label{q3:varsum}
\Var(\hat\beta_1)=\frac{1}{\mathrm{SST}_x^2}\Big[\sum_t x_t^2\Var(u_t)+2\sum_{t<s}x_tx_s\Cov(u_t,u_s)\Big].
\end{equation}
Substituting \eqref{q3:cov} into \eqref{q3:varsum}:
\begin{equation}\label{q3:true}
\Var(\hat\beta_1)=\frac{\sigma^2}{\mathrm{SST}_x}
+\frac{2\sigma^2}{\mathrm{SST}_x^2}\sum_{t=1}^{n-1}\sum_{j=1}^{n-t}\rho^{j}x_tx_{t+j},
\end{equation}
which is Wooldridge's equation (12.4).

\passo{What the usual formula estimates}{The usual OLS variance is}
\begin{equation}\label{q3:usual}
\widehat{\Var}_{\mathrm{OLS}}(\hat\beta_1)=\frac{\hat\sigma^2}{\mathrm{SST}_x},
\end{equation}
an estimate of the first term of \eqref{q3:true} only. The usual standard error is therefore biased
by the omitted second term.

\passo{Why it is usually too small}{In economic data $\rho>0$ and the regressors are positively
autocorrelated (persistent, trending), so $x_tx_{t+j}>0$ on average and every $\rho^j>0$: the double
sum in \eqref{q3:true} is positive and \eqref{q3:usual} understates the true variance. A large-$n$
illustration: if $x_t$ is AR(1) with autocorrelation $\lambda$, then
$\sum_{t}\sum_{j}\rho^jx_tx_{t+j}\approx\mathrm{SST}_x\sum_{j\ge1}(\rho\lambda)^j$, and substituting into
\eqref{q3:true}:}
\begin{equation}\label{q3:ratio}
\Var(\hat\beta_1)\approx\frac{\sigma^2}{\mathrm{SST}_x}\Big[1+\frac{2\rho\lambda}{1-\rho\lambda}\Big]
=\frac{\sigma^2}{\mathrm{SST}_x}\cdot\frac{1+\rho\lambda}{1-\rho\lambda}.
\end{equation}
With $\rho=\lambda=0.5$, \eqref{q3:ratio} gives the ratio $(1+0.25)/(1-0.25)=1.25/0.75=1.6667$: the true
standard error is $\sqrt{1.6667}=1.2910$ times the usual one.

\passo{Not always}{The sign of the second term in \eqref{q3:true} is the sign of
$\sum\rho^jx_tx_{t+j}$. If $\rho<0$ (with persistent $x$), or if $\rho>0$ but $x_t$ alternates in sign
(negatively autocorrelated), the cross terms are negative on net and the usual standard error is too
large. In \eqref{q3:ratio}, $\rho=-0.5$, $\lambda=0.5$ gives $(1-0.25)/(1+0.25)=0.75/1.25=0.6000$, so the
true standard error is only $\sqrt{0.6}=0.7746$ times the usual one. If $x_t$ is serially uncorrelated
($\lambda=0$) the bias vanishes in large samples even with $\rho\neq0$. So the answer to the second
question is no.}

\intuicao{Positively correlated errors make neighbouring observations repeat each other, so the
sample carries less independent information than $n$ suggests; the usual formula counts every
observation as new. When regressor and error both drift slowly, a run of positive errors coincides
with a run of large $x$, which tilts the fitted slope more than independent errors would.}
\refbib{Wooldridge (2016), §12-1b (efficiency and inference, pp.~373--374).}
\end{sol}

% =============================================================================
\bloco{3}{Greenberg, \textit{Introduction to Bayesian Econometrics}}{G 2.2; G 7.3}

\questao{G 2.2 --- gamma--Poisson conjugacy.}
\fonte{Greenberg, E. (2008), \textit{Introduction to Bayesian Econometrics}, 1st ed., Cambridge
University Press, §2.5 Exercises, Exercise~2.2, p.~19 --- reproduced verbatim.}

2.2 The Poisson distribution has probability mass function
\[ p(y_i\mid\theta)=\frac{\theta^{y_i}e^{-\theta}}{y_i!},\qquad \theta>0,\ y_i=0,1,\dots, \]
and let $y_1,\dots,y_n$ be a random sample from this distribution.
\begin{enumerate}[label=\textbf{(\alph*)},leftmargin=2.4em]
\item Show that the gamma distribution $G(\alpha,\beta)$ is a conjugate prior distribution for the
Poisson distribution.
\item Show that $\bar y$ is the MLE for $\theta$.
\item Write the mean of the posterior distribution as a weighted average of the mean of the prior
distribution and the MLE.
\item What happens to the weight on the prior mean as $n$ becomes large?
\end{enumerate}
\resolespaco
\begin{sol}
\solhead
\begin{notacao}
\item $y_i\in\{0,1,\dots\}$: counts; $\theta>0$: Poisson rate; $s=\sum_{i=1}^ny_i$, $\bar y=s/n$.
\item $G(\alpha,\beta)$, $\alpha,\beta>0$: gamma with shape $\alpha$ and rate $\beta$ (Greenberg's
  convention, App.~A), $\pi(\theta)=\frac{\beta^\alpha}{\Gamma(\alpha)}\theta^{\alpha-1}e^{-\beta\theta}$,
  mean $\alpha/\beta$.
\item $\ell(\theta)$: log-likelihood; $\hat\theta$: MLE; $w$: weight on the prior mean.
\end{notacao}

\passo{(a) Likelihood}{Independence gives}
\begin{equation}\label{q4:lik}
p(y\mid\theta)=\prod_{i=1}^n\frac{\theta^{y_i}e^{-\theta}}{y_i!}=\frac{\theta^{s}e^{-n\theta}}{\prod_iy_i!}.
\end{equation}

\passo{(a) Posterior}{Bayes' theorem, $\pi(\theta\mid y)\propto p(y\mid\theta)\pi(\theta)$, with
\eqref{q4:lik} and the gamma prior, dropping factors free of $\theta$:}
\begin{equation}\label{q4:post0}
\pi(\theta\mid y)\propto\big(\theta^{s}e^{-n\theta}\big)\big(\theta^{\alpha-1}e^{-\beta\theta}\big).
\end{equation}
Adding exponents in \eqref{q4:post0}:
\begin{equation}\label{q4:post}
\pi(\theta\mid y)\propto\theta^{\alpha+s-1}e^{-(\beta+n)\theta},
\end{equation}
the kernel of $G(\alpha+s,\ \beta+n)$. Prior and posterior are both gamma, so the gamma family is
conjugate for the Poisson likelihood.

\passo{(b) MLE}{Taking logs of \eqref{q4:lik}:}
\begin{equation}\label{q4:ll}
\ell(\theta)=s\log\theta-n\theta-\sum_i\log y_i!.
\end{equation}
Differentiating \eqref{q4:ll} and setting the score to zero:
\begin{equation}\label{q4:score}
\ell'(\theta)=\frac{s}{\theta}-n=0\;\Longrightarrow\;\hat\theta=\frac{s}{n}=\bar y .
\end{equation}
The second derivative $\ell''(\theta)=-s/\theta^2<0$ when $s>0$, so \eqref{q4:score} is a maximum. (If
$s=0$, $\ell(\theta)=-n\theta$ decreases and the supremum is at the boundary $0=\bar y$.)

\passo{(c) Posterior mean as a weighted average}{The mean of the gamma \eqref{q4:post} is}
\begin{equation}\label{q4:pm}
\E(\theta\mid y)=\frac{\alpha+s}{\beta+n}.
\end{equation}
Substituting $\alpha=\beta(\alpha/\beta)$ and $s=n\bar y$ into \eqref{q4:pm}:
\[
\E(\theta\mid y)=\frac{\beta\,(\alpha/\beta)+n\,(\bar y)}{\beta+n}.
\]
Splitting the fraction:
\begin{equation}\label{q4:wavg}
\E(\theta\mid y)=\underbrace{\frac{\beta}{\beta+n}}_{w}\cdot\frac{\alpha}{\beta}
+\underbrace{\frac{n}{\beta+n}}_{1-w}\cdot\bar y .
\end{equation}

\passo{(d) Large $n$}{From \eqref{q4:wavg},}
\begin{equation}\label{q4:lim}
w=\frac{\beta}{\beta+n}\xrightarrow[n\to\infty]{}0,
\end{equation}
so the posterior mean converges to the MLE $\bar y$; the posterior variance $(\alpha+s)/(\beta+n)^2$
also goes to zero at rate $1/n$.

\passo{Numerical illustration}{Coal-mining disasters 1851--1891 ($n=41$, $s=127$,
$\bar y=3.097561$) with the prior of G~7.3, $\alpha=0.5$, $\beta=1$. By \eqref{q4:lim},
$w=1/(1+41)=0.023810$; by \eqref{q4:wavg},}
\[
\E(\theta\mid y)=0.023810\,(0.5)+0.976190\,(3.097561)=0.011905+3.023810=3.035714,
\]
which equals \eqref{q4:pm}: $(0.5+127)/(1+41)=127.5/42=3.035714$.

\intuicao{The prior acts like $\beta$ extra periods of data containing $\alpha$ events in total.
After $n$ real periods it weighs $\beta/(\beta+n)$; with 41 years of data a prior worth one year
moves the estimate only from $3.098$ to $3.036$.}
\refbib{Greenberg (2008), §2.2 (pp.~12--18) and App.~A (gamma distribution); the Class~1 deck
\textit{first Bayesian example} (Poisson with an exponential prior, the case $\alpha=1$).}
\end{sol}

\questao{G 7.3 --- Gibbs sampler for the coal-mining switch point (computer exercise).}
\fonte{Greenberg, E. (2008), \textit{Introduction to Bayesian Econometrics}, 1st ed., Cambridge
University Press, §7.5 Exercises, Exercise~7.3, p.~106 --- reproduced verbatim. Equations (7.3) and
(7.4), to which the exercise refers, are reproduced verbatim from §7.1.1, pp.~93--94.}

\emph{Equations referred to (Greenberg 2008, pp.~93--94):}
\[
p(y_i)=\begin{cases}\dfrac{e^{-\theta_1}\theta_1^{y_i}}{y_i!}, & \text{for } i=1,\dots,k,\\[8pt]
\dfrac{e^{-\theta_2}\theta_2^{y_i}}{y_i!}, & \text{for } i=k+1,\dots,n,\end{cases}\tag{7.3}
\]
``where $y_i=0,1,\dots,$ and the switch point $k$ is unknown. The specification is completed by
assigning the conditionally conjugate priors,
\[
\theta_1\sim G(\alpha_{10},\beta_{10}),\quad \theta_2\sim G(\alpha_{20},\beta_{20}),\quad
\pi(k=j)=1/n,\quad j=1,\dots,n.\text{''}\tag{7.4}
\]

7.3 Construct a Gibbs algorithm to analyze the Poisson model with unknown switch point. Given the
specification in Equations (7.3) and (7.4), show that
\[
\pi(\theta_1,\theta_2,k\mid y)\propto\theta_1^{\alpha_{10}-1}e^{-\beta_{10}\theta_1}\theta_2^{\alpha_{20}-1}e^{-\beta_{20}\theta_2}
\prod_{i=1}^{k}e^{-\theta_1}\theta_1^{y_i}\prod_{i=k+1}^{n}e^{-\theta_2}\theta_2^{y_i},
\]
and verify that
\begin{gather*}
\theta_1\mid y,k\sim G\Big(\alpha_{10}+\sum_{1}^{k}y_i,\ \beta_{10}+k\Big),\\
\theta_2\mid k,y\sim G\Big(\alpha_{20}+\sum_{k+1}^{n}y_i,\ \beta_{20}+n-k\Big),\\
\pi(k\mid y,\theta_1,\theta_2)=\frac{e^{k(\theta_2-\theta_1)}(\theta_1/\theta_2)^{\sum_1^ky_i}}
{\sum_{k=1}^{n}e^{k(\theta_2-\theta_1)}\,(\theta_1/\theta_2)^{\sum_1^ky_i}},\qquad k=1,\dots,n.
\end{gather*}
Apply your algorithm to the mining disaster data analyzed by Carlin et al.\ (1992a); the data may be
found in the article. Your analysis should generate marginal posterior distributions for
$\theta_1$, $\theta_2$, $p(k)$, and the marginal likelihood computed by the method of Chib (1995).
For simplicity, set $\alpha_{10}=\alpha_{20}=0.5$ and $\beta_{10}=\beta_{20}=1.0$. (Curiously,
although these data have been analyzed from the Bayesian viewpoint in several papers, none of them
present historical material that might shed light on when changes in parameters took place. For
example, changes in government regulation, technology, or the location of mines might be expected
to affect the probability of a disaster. A more general version of the model permits multiple
switch points if necessary.)
\resolespaco
\begin{sol}
\solhead
\begin{notacao}
\item $y_i$: number of British coal-mining disasters in year $i$, $i=1$ (1851) to $n=112$ (1962);
  the series of Carlin, Gelfand \& Smith (1992) is Jarrett (1979), total $S=\sum_iy_i=191$.
\item $k\in\{1,\dots,n\}$: switch point, the last year of regime 1 (year $1850+k$); $k=n$ means no change.
\item $\theta_1,\theta_2>0$: disaster rates before and after the switch;
  $S_k=\sum_{i=1}^ky_i$ (so $\sum_{k+1}^ny_i=S-S_k$).
\item $G(a,b)$: gamma with shape $a$, rate $b$; $\alpha_{10}=\alpha_{20}=\alpha=0.5$,
  $\beta_{10}=\beta_{20}=\beta=1$.
\item $g=1,\dots,G$: Gibbs iteration; $\theta^*=(\theta_1^*,\theta_2^*)$: point where Chib's
  identity is evaluated; $m(y)$: marginal likelihood.
\end{notacao}

\passo{Joint posterior}{By (7.3) and independence, the likelihood is}
\begin{equation}\label{q5:lik}
p(y\mid\theta_1,\theta_2,k)=\prod_{i=1}^k\frac{e^{-\theta_1}\theta_1^{y_i}}{y_i!}
\prod_{i=k+1}^n\frac{e^{-\theta_2}\theta_2^{y_i}}{y_i!},
\end{equation}
and by (7.4) the prior is
\begin{equation}\label{q5:prior}
\pi(\theta_1,\theta_2,k)=\frac{\beta_{10}^{\alpha_{10}}}{\Gamma(\alpha_{10})}\theta_1^{\alpha_{10}-1}e^{-\beta_{10}\theta_1}
\cdot\frac{\beta_{20}^{\alpha_{20}}}{\Gamma(\alpha_{20})}\theta_2^{\alpha_{20}-1}e^{-\beta_{20}\theta_2}\cdot\frac1n .
\end{equation}
Multiplying \eqref{q5:lik} by \eqref{q5:prior} and dropping the factors free of
$(\theta_1,\theta_2,k)$ (the $y_i!$, the gamma constants, $1/n$) gives the expression in the statement:
\begin{equation}\label{q5:joint}
\pi(\theta_1,\theta_2,k\mid y)\propto\theta_1^{\alpha_{10}-1}e^{-\beta_{10}\theta_1}\theta_2^{\alpha_{20}-1}e^{-\beta_{20}\theta_2}
\prod_{i=1}^{k}e^{-\theta_1}\theta_1^{y_i}\prod_{i=k+1}^{n}e^{-\theta_2}\theta_2^{y_i}.
\end{equation}

\passo{Full conditional of $\theta_1$}{Keeping in \eqref{q5:joint} only the factors with $\theta_1$:}
\[
\pi(\theta_1\mid y,k,\theta_2)\propto\theta_1^{\alpha_{10}-1}e^{-\beta_{10}\theta_1}\big(e^{-k\theta_1}\theta_1^{S_k}\big).
\]
Adding exponents:
\begin{equation}\label{q5:t1}
\pi(\theta_1\mid y,k)\propto\theta_1^{\alpha_{10}+S_k-1}e^{-(\beta_{10}+k)\theta_1}
\;\Longrightarrow\;\theta_1\mid y,k\sim G(\alpha_{10}+S_k,\ \beta_{10}+k).
\end{equation}

\passo{Full conditional of $\theta_2$}{In the same way, the $\theta_2$ factors of \eqref{q5:joint}
are}
\[
\pi(\theta_2\mid y,k,\theta_1)\propto\theta_2^{\alpha_{20}-1}e^{-\beta_{20}\theta_2}\big(e^{-(n-k)\theta_2}\theta_2^{S-S_k}\big),
\]
so, adding exponents,
\begin{equation}\label{q5:t2}
\theta_2\mid y,k\sim G(\alpha_{20}+S-S_k,\ \beta_{20}+n-k).
\end{equation}
Neither \eqref{q5:t1} nor \eqref{q5:t2} involves the other rate: given $k$, $\theta_1$ and $\theta_2$
are independent.

\passo{Full conditional of $k$}{The factors of \eqref{q5:joint} that change with $k$ are the two
products:}
\begin{equation}\label{q5:k0}
\pi(k\mid y,\theta_1,\theta_2)\propto e^{-k\theta_1}\theta_1^{S_k}\,e^{-(n-k)\theta_2}\theta_2^{S-S_k}.
\end{equation}
Dividing \eqref{q5:k0} by the $k$-free factor $e^{-n\theta_2}\theta_2^{S}$:
\begin{equation}\label{q5:k}
\pi(k\mid y,\theta_1,\theta_2)\propto e^{k(\theta_2-\theta_1)}\Big(\frac{\theta_1}{\theta_2}\Big)^{S_k},
\qquad k=1,\dots,n,
\end{equation}
and normalising over $k=1,\dots,n$ gives the formula in the statement.

\passo{Algorithm}{Start at $k^{(0)}=56$. For $g=1,\dots,11{,}000$: draw $\theta_1^{(g)}$ from
\eqref{q5:t1} and $\theta_2^{(g)}$ from \eqref{q5:t2} given $k^{(g-1)}$; compute the $n$ weights
\eqref{q5:k} on the log scale, normalise, and draw $k^{(g)}$. Discard 1{,}000 burn-in draws and keep
$G=10{,}000$ (NumPy seed 2026).}

\passo{Exact benchmark}{Because the priors are conjugate given $k$, $\theta_1,\theta_2$ integrate out
of \eqref{q5:lik}$\times$\eqref{q5:prior} in closed form:}
\begin{equation}\label{q5:pyk}
p(y\mid k)=\frac{1}{\prod_iy_i!}\cdot\frac{\beta^\alpha\,\Gamma(\alpha+S_k)}{\Gamma(\alpha)\,(\beta+k)^{\alpha+S_k}}
\cdot\frac{\beta^\alpha\,\Gamma(\alpha+S-S_k)}{\Gamma(\alpha)\,(\beta+n-k)^{\alpha+S-S_k}} ,
\end{equation}
so the exact marginal likelihood and posterior of $k$ are
\begin{equation}\label{q5:exact}
m(y)=\sum_{k=1}^n\frac1n\,p(y\mid k),\qquad p(k\mid y)=\frac{p(y\mid k)/n}{m(y)},
\end{equation}
giving $\log m(y)=-179.4590$, $\E(\theta_1\mid y)=3.0500$, $\E(\theta_2\mid y)=0.9147$. These check the
sampler.

\passo{Marginal posteriors}{Gibbs output (exact values from \eqref{q5:exact} for $k$):}
\begin{center}\small
\begin{tabular}{lccc}
\toprule
 & Posterior mean & Posterior sd & 95\% interval \\
\midrule
$\theta_1$ (before) & 3.0506 & 0.2809 & $[2.530,\ 3.635]$\\
$\theta_2$ (after)  & 0.9166 & 0.1163 & $[0.701,\ 1.156]$\\
\bottomrule
\end{tabular}
\qquad
\begin{tabular}{cccc}
\toprule
$k$ & year & $p(k\mid y)$ Gibbs & exact \\
\midrule
41 & 1891 & 0.2471 & 0.2482\\
40 & 1890 & 0.1849 & 0.1848\\
39 & 1889 & 0.1433 & 0.1414\\
42 & 1892 & 0.0987 & 0.1010\\
37 & 1887 & 0.0881 & 0.0898\\
36 & 1886 & 0.0742 & 0.0745\\
\bottomrule
\end{tabular}
\end{center}
The posterior of $k$ has mean $40.14$ and mode $k=41$ (1891); $P(36\le k\le46\mid y)=0.98$, i.e.\ the
switch happened between 1886 and 1896 with high probability. Disasters fell from about 3.05 to
about 0.92 per year.

\passo{Marginal likelihood by Chib (1995)}{For any $\theta^*$, Bayes' theorem rearranged gives the
identity}
\begin{equation}\label{q5:chib}
\log m(y)=\log p(y\mid\theta^*)+\log\pi(\theta^*)-\log\pi(\theta^*\mid y).
\end{equation}
Here $k$ is a latent parameter, so the likelihood term sums it out with its prior:
\begin{equation}\label{q5:lstar}
p(y\mid\theta^*)=\sum_{k=1}^n\frac1n\,p(y\mid\theta_1^*,\theta_2^*,k).
\end{equation}
For the ordinate, conditional independence of $\theta_1,\theta_2$ given $k$ (from \eqref{q5:t1} and
\eqref{q5:t2}) gives the Rao--Blackwell estimate
\begin{equation}\label{q5:ord}
\hat\pi(\theta^*\mid y)=\frac1G\sum_{g=1}^G G\big(\theta_1^*\mid\alpha+S_{k^{(g)}},\beta+k^{(g)}\big)\,
G\big(\theta_2^*\mid\alpha+S-S_{k^{(g)}},\beta+n-k^{(g)}\big),
\end{equation}
with $G(\cdot\mid a,b)$ the gamma density. At the posterior means $\theta^*=(3.0506,\ 0.9166)$:
$\log p(y\mid\theta^*)=-172.2563$ from \eqref{q5:lstar}, $\log\pi(\theta^*)=-5.6261$ from (7.4), and
$\log\hat\pi(\theta^*\mid y)=1.5765$ from \eqref{q5:ord}. Substituting into \eqref{q5:chib}:
\[
\log\hat m(y)=(-172.2563)+(-5.6261)-(1.5765)=-179.4589 ,
\]
against the exact $-179.4590$ from \eqref{q5:exact}.

\passo{Python}{Core of the sampler (complete code with assertions in the check script):}
{\small
\begin{verbatim}
cs, ks = np.cumsum(y), np.arange(1, n + 1)        # S_k and k = 1..n
rng, k = np.random.default_rng(2026), n // 2
for g in range(11000):
    t1 = rng.gamma(a0 + cs[k-1], 1 / (b0 + k))                 # (5.4) theta1 | k
    t2 = rng.gamma(a0 + S - cs[k-1], 1 / (b0 + n - k))         # (5.5) theta2 | k
    lp = ks * (t2 - t1) + cs * np.log(t1 / t2)                 # (5.7) log pi(k | ...)
    p = np.exp(lp - lp.max()); k = rng.choice(ks, p=p / p.sum())
# output: theta1 3.0506 (0.2809), theta2 0.9166 (0.1163), mode k = 41 (1891)
# Chib: -172.2563 + (-5.6261) - 1.5765 = -179.4589 ; exact -179.4590
\end{verbatim}
}

\intuicao{The data choose a sharp break: about three disasters a year until around 1890, under one
afterwards, with the switch year pinned down within a few years. The Gibbs sampler works because
each block is easy given the others: rates are conjugate gammas given $k$, and $k$ is a discrete
draw over 112 years given the rates. On the historical question Greenberg raises, the posterior
mass near 1887--1891 sits close to the Coal Mines Regulation Act of 1887, a candidate explanation
that this model does not test. Chib's identity recovers the exact marginal likelihood to the fourth
decimal, which is what makes it usable for comparing this model with, say, a no-break Poisson in
Classes 3+4.}
\refbib{Greenberg (2008), §7.1.1 (pp.~91--95) and §7.4 (Chib 1995); Class~1 deck
\textit{Poisson with breakpoint} and \texttt{Notas/L01/L01-03-coal-mining-change-point} (the profile
likelihood also peaks at 1891).}
\end{sol}

% =============================================================================
\bloco{4}{James, Witten, Hastie, Tibshirani \& Taylor, \textit{ISLP}}{I 5.2; I 5.9}

\questao{I 5.2 --- probability that an observation is in a bootstrap sample.}
\fonte{James, G., Witten, D., Hastie, T., Tibshirani, R. and Taylor, J. (2023), \textit{An
Introduction to Statistical Learning with Applications in Python}, Springer, §5.4 Exercises
(Conceptual), Exercise~2, pp.~224--225 --- reproduced verbatim (the margin index entry
``\texttt{np.empty()}'' is omitted).}

2. We will now derive the probability that a given observation is part of a bootstrap sample.
Suppose that we obtain a bootstrap sample from a set of $n$ observations.
\begin{enumerate}[label=(\alph*),leftmargin=2.4em]
\item What is the probability that the first bootstrap observation is \emph{not} the $j$th
observation from the original sample? Justify your answer.
\item What is the probability that the second bootstrap observation is \emph{not} the $j$th
observation from the original sample?
\item Argue that the probability that the $j$th observation is \emph{not} in the bootstrap sample
is $(1-1/n)^n$.
\item When $n=5$, what is the probability that the $j$th observation is in the bootstrap sample?
\item When $n=100$, what is the probability that the $j$th observation is in the bootstrap sample?
\item When $n=10{,}000$, what is the probability that the $j$th observation is in the bootstrap
sample?
\item Create a plot that displays, for each integer value of $n$ from 1 to 100,000, the probability
that the $j$th observation is in the bootstrap sample. Comment on what you observe.
\item We will now investigate numerically the probability that a bootstrap sample of size $n=100$
contains the $j$th observation. Here $j=4$. We first create an array \texttt{store} with values that
will subsequently be overwritten using the function \texttt{np.empty()}. We then repeatedly create
bootstrap samples, and each time we record whether or not the fifth observation is contained in the
bootstrap sample.
{\small
\begin{verbatim}
rng = np.random.default_rng(10)
store = np.empty(10000)
for i in range(10000):
    store[i] = np.sum(rng.choice(100, replace=True) == 4)
          > 0
np.mean(store)
\end{verbatim}
}
Comment on the results obtained.
\end{enumerate}
\resolespaco
\begin{sol}
\solhead
\begin{notacao}
\item $n$: sample size; $j\in\{1,\dots,n\}$: a fixed original observation.
\item $Z_1,\dots,Z_n$: indices of the bootstrap draws, i.i.d.\ uniform on $\{1,\dots,n\}$ (sampling
  with replacement).
\item $p_n=P(j\text{ is in the bootstrap sample})$.
\end{notacao}

\passo{(a)}{Each draw picks one of the $n$ original observations with equal probability $1/n$, so}
\begin{equation}\label{q6:a}
P(Z_1\neq j)=1-\frac1n .
\end{equation}

\passo{(b)}{Sampling is with replacement, so $Z_2$ has the same distribution as $Z_1$ and is
independent of it:}
\begin{equation}\label{q6:b}
P(Z_2\neq j)=1-\frac1n .
\end{equation}

\passo{(c)}{$j$ is absent only if every draw misses it; by independence and \eqref{q6:a}--\eqref{q6:b}
applied to all $n$ draws:}
\begin{equation}\label{q6:c}
P(j\notin\text{sample})=\prod_{m=1}^nP(Z_m\neq j)=\Big(1-\frac1n\Big)^n .
\end{equation}

\passo{(d)--(f)}{Taking the complement of \eqref{q6:c}:}
\begin{equation}\label{q6:p}
p_n=1-\Big(1-\frac1n\Big)^n .
\end{equation}
Substituting into \eqref{q6:p}:
\begin{align*}
n=5:&\quad p_5=1-(4/5)^5=1-0.32768=0.67232,\\
n=100:&\quad p_{100}=1-(0.99)^{100}=1-0.366032=0.633968,\\
n=10{,}000:&\quad p_{10000}=1-(0.9999)^{10000}=1-0.367861=0.632139.
\end{align*}

\passo{(g) Plot and limit}{Since $n\log(1-1/n)=n\big(-\tfrac1n-\tfrac1{2n^2}-\dots\big)\to-1$,}
\begin{equation}\label{q6:lim}
\lim_{n\to\infty}p_n=1-e^{-1}=0.632121 .
\end{equation}
\begin{center}
\begin{tikzpicture}
\begin{semilogxaxis}[width=0.72\linewidth,height=5cm,xlabel={$n$ (log scale)},
  ylabel={$p_n$},ymin=0.6,ymax=1.02,xmin=1,xmax=100000,grid=major,
  every axis plot/.append style={thick}]
\addplot[blue!60!black,mark=*,mark size=1.3pt] coordinates {(1,1.0000) (2,0.7500) (3,0.7037)
  (5,0.6723) (10,0.6513) (20,0.6415) (50,0.6358) (100,0.6340) (1000,0.6323) (10000,0.6321)
  (100000,0.6321)};
\addplot[red!70!black,dashed,domain=1:100000] {0.632121};
\end{semilogxaxis}
\end{tikzpicture}
\end{center}
$p_n$ falls from 1 at $n=1$ and flattens almost immediately at $1-e^{-1}\approx0.632$ (dashed): for
$n\ge100$ the curve is indistinguishable from the limit.

\passo{(h) The printed code}{\texttt{rng.choice(100, replace=True)} without \texttt{size} returns a
\emph{single} integer, so each iteration checks whether one draw equals 4, an event of probability
$1/100$. Run as printed, \texttt{np.mean(store)} returns $0.0089$, not a bootstrap probability. The
intended code draws a whole bootstrap sample, \texttt{rng.choice(100, size=100, replace=True)},
and then returns $0.6362$, close to $p_{100}=0.633968$ from (e) (Monte Carlo standard error
$\sqrt{0.634\cdot0.366/10000}\approx0.005$). Index 4 is the fifth observation because Python
counts from 0, which is why the text says ``fifth''.}

\intuicao{A bootstrap sample of size $n$ contains only about 63.2\% of the distinct original
observations; the other 36.8\% are ``out of bag''. This is why bootstrap estimates of test error
that reuse the training points are optimistic, and why out-of-bag points can serve as a free
validation set.}
\refbib{ISLP (2023), §5.2 (the bootstrap, pp.~212--215).}
\end{sol}

\questao{I 5.9 --- bootstrap standard errors for Boston \texttt{medv} (computer exercise).}
\fonte{James, G., Witten, D., Hastie, T., Tibshirani, R. and Taylor, J. (2023), \textit{An
Introduction to Statistical Learning with Applications in Python}, Springer, §5.4 Exercises
(Applied), Exercise~9, pp.~227--228 --- reproduced verbatim (the margin index entry
``\texttt{np.percentile()}'' is omitted).}

9. We will now consider the \texttt{Boston} housing data set, from the \texttt{ISLP} library.
\begin{enumerate}[label=(\alph*),leftmargin=2.4em]
\item Based on this data set, provide an estimate for the population mean of \texttt{medv}. Call
this estimate $\hat\mu$.
\item Provide an estimate of the standard error of $\hat\mu$. Interpret this result.\\
\textit{Hint: We can compute the standard error of the sample mean by dividing the sample standard
deviation by the square root of the number of observations.}
\item Now estimate the standard error of $\hat\mu$ using the bootstrap. How does this compare to
your answer from (b)?
\item Based on your bootstrap estimate from (c), provide a 95\,\% confidence interval for the mean
of \texttt{medv}. Compare it to the results obtained by using \texttt{Boston['medv'].std()} and the
two standard error rule (3.9).\\
\textit{Hint: You can approximate a 95\,\% confidence interval using the formula
$[\hat\mu-2\mathrm{SE}(\hat\mu),\hat\mu+2\mathrm{SE}(\hat\mu)]$.}
\item Based on this data set, provide an estimate, $\hat\mu_{med}$, for the median value of
\texttt{medv} in the population.
\item We now would like to estimate the standard error of $\hat\mu_{med}$. Unfortunately, there is
no simple formula for computing the standard error of the median. Instead, estimate the standard
error of the median using the bootstrap. Comment on your findings.
\item Based on this data set, provide an estimate for the tenth percentile of \texttt{medv} in
Boston census tracts. Call this quantity $\hat\mu_{0.1}$. (You can use the \texttt{np.percentile()}
function.)
\item Use the bootstrap to estimate the standard error of $\hat\mu_{0.1}$. Comment on your findings.
\end{enumerate}
\resolespaco
\begin{sol}
\solhead
\begin{notacao}
\item $x_i$: \texttt{medv} of tract $i$, median value of owner-occupied homes in thousands of US
  dollars; $N=506$ tracts (file \texttt{Listas/data/Boston.csv}, from the ISLP repository).
\item $\hat\mu$: sample mean; $s$: sample standard deviation (divisor $N-1$); $\mathrm{SE}$: standard
  error.
\item $B=10{,}000$ bootstrap samples $x^{*b}$ of size $N$ drawn with replacement (NumPy seed 1);
  $\hat\mu^{*b}$, $\hat\mu^{*b}_{med}$, $\hat\mu^{*b}_{0.1}$: the statistic on sample $b$.
\end{notacao}

\passo{(a) Mean}{}
\begin{equation}\label{q7:mean}
\hat\mu=\frac1N\sum_{i=1}^Nx_i=22.5328 .
\end{equation}

\passo{(b) Formula standard error}{}
\begin{equation}\label{q7:se}
\mathrm{SE}(\hat\mu)=\frac{s}{\sqrt N}.
\end{equation}
Substituting into \eqref{q7:se}:
\[
\mathrm{SE}(\hat\mu)=\frac{9.1971}{\sqrt{506}}=\frac{9.1971}{22.4944}=0.4089 .
\]
If we could draw many samples of 506 tracts, their means would scatter around the population mean
with a standard deviation of about 0.41 (roughly 409 US dollars).

\passo{(c) Bootstrap standard error}{}
\begin{equation}\label{q7:seb}
\mathrm{SE}_B(\hat\mu)=\sqrt{\frac1{B-1}\sum_{b=1}^B\Big(\hat\mu^{*b}-\frac1B\sum_{b'=1}^B\hat\mu^{*b'}\Big)^2}=0.4098 .
\end{equation}
It agrees with \eqref{q7:se} to within $0.0009$, the size of the Monte Carlo error of $B=10{,}000$
replicates: for the mean, the bootstrap reproduces the textbook formula.

\passo{(d) 95\% intervals}{Using the two-standard-error rule with \eqref{q7:seb}:}
\begin{equation}\label{q7:ci}
[\hat\mu-2\,\mathrm{SE},\ \hat\mu+2\,\mathrm{SE}].
\end{equation}
Substituting into \eqref{q7:ci}:
\begin{align*}
&[\,22.5328-2(0.4098),\ 22.5328+2(0.4098)\,]\\
&=[\,22.5328-0.8196,\ 22.5328+0.8196\,]=[\,21.7132,\ 23.3524\,].
\end{align*}
With \texttt{Boston['medv'].std()} $=s$ in \eqref{q7:se}: $[\,22.5328-0.8177,\ 22.5328+0.8177\,]=[\,21.7151,\ 23.3505\,]$.
The two intervals differ only in the third decimal.

\passo{(e) Median}{With $N$ even, $\hat\mu_{med}$ is the average of the 253rd and 254th ordered values:
$\hat\mu_{med}=21.2$.}

\passo{(f) Bootstrap SE of the median}{Applying \eqref{q7:seb} to $\hat\mu^{*b}_{med}$:
$\mathrm{SE}_B(\hat\mu_{med})=0.3796$. The median is estimated about as precisely as the mean, even a
little more: \texttt{medv} is right-skewed and capped at 50, and the median ignores the tail
observations that inflate the variance of the mean.}

\passo{(g) Tenth percentile}{\texttt{np.percentile(medv, 10)} (linear interpolation) gives
$\hat\mu_{0.1}=12.75$.}

\passo{(h) Bootstrap SE of the tenth percentile}{Applying \eqref{q7:seb} to $\hat\mu^{*b}_{0.1}$:
$\mathrm{SE}_B(\hat\mu_{0.1})=0.5080$, larger than for the median. Fewer observations lie near the 10th
percentile, so it moves more from sample to sample; asymptotically
$\mathrm{SE}(\hat q_p)\approx\sqrt{p(1-p)/N}/f(q_p)$, which grows where the density $f$ is low.}

\passo{Python}{}
{\small
\begin{verbatim}
medv = pd.read_csv("Listas/data/Boston.csv")["medv"].to_numpy()   # N = 506
mu, se = medv.mean(), medv.std(ddof=1) / np.sqrt(len(medv))       # 22.5328, 0.4089
rng = np.random.default_rng(1)
boot = medv[rng.integers(0, 506, size=(10000, 506))]              # B = 10,000 samples
boot.mean(1).std(ddof=1)                     # 0.4098
np.median(boot, 1).std(ddof=1)               # 0.3796   (median 21.2)
np.percentile(boot, 10, axis=1).std(ddof=1)  # 0.5080   (10th percentile 12.75)
\end{verbatim}
}

\intuicao{The bootstrap replaces a formula by resampling: for the mean it reproduces $s/\sqrt N$, and
for the median and percentiles, which have no simple formula, it is the practical tool. The SEs
also show that tail quantiles are noisier than central ones.}
\refbib{ISLP (2023), §5.2 (the bootstrap, pp.~212--215) and §3.1.2 (the two standard error rule (3.9)).}
\end{sol}

% =============================================================================
\bloco{5}{Bolfarine \& Sandoval, \textit{Introdução à Inferência Estatística}}{B 1.4, 1.9, 1.12 $\to$ 2.2, 2.4, 2.10 $\to$ 3.6, 3.13, 3.17}

\questao{B 1.4 --- MSE of $S^2$ and $\hat\sigma^2$.}
\fonte{Bolfarine, H. and Sandoval, M.~C. (2010), \textit{Introdução à Inferência Estatística},
§1.4 Exercícios, Exercício~1.4, p.~13 --- reproduced verbatim. The expressions it refers to are
reproduced verbatim from Exemplo~1.3.6, p.~9.}

\begin{otherlanguage}{brazil}
\textbf{1.4.} Verifique a validade das expressões (1.3.10) e (1.3.11).

\smallskip
\emph{Expressões referidas (Exemplo 1.3.6, p.~9):} ``Sejam $X_1,\dots,X_n$ uma amostra aleatória da
variável aleatória $X\sim N(\mu,\sigma^2)$. Conforme visto no Exemplo 1.3.3, $\hat\sigma^2$ é um
estimador viciado para $\sigma^2$. De (1.3.3) segue que
\[ S^2=\frac{n}{n-1}\hat\sigma^2=\frac1{n-1}\sum_{i=1}^n(X_i-\overline X)^2 \]
é um estimador não viciado para $\sigma^2$. Por outro lado, temos (ver Exercício 1.4) que
\[ EQM[S^2]=Var[S^2]=\frac{2\sigma^4}{n-1}, \tag{1.3.10}\]
e que
\[ EQM[\hat\sigma^2]=\frac{2\sigma^4}{(n-1)}\left[1-\frac{(3n-1)}{2n^2}\right].\text{''} \tag{1.3.11}\]
\end{otherlanguage}
\resolespaco
\begin{sol}
\solhead
\begin{notacao}
\item $X_1,\dots,X_n$ i.i.d.\ $N(\mu,\sigma^2)$, $n\ge2$; $\overline X$: sample mean.
\item $\hat\sigma^2=\frac1n\sum_i(X_i-\overline X)^2$ (MLE, Exemplo 1.3.3); $S^2=\frac1{n-1}\sum_i(X_i-\overline X)^2$.
\item $EQM[T]=\E[(T-\sigma^2)^2]$: mean squared error; $\chi^2_\nu$: chi-square with $\nu$ degrees of freedom.
\end{notacao}

\passo{Sampling distribution}{For a normal sample (Student--Fisher theorem),}
\begin{equation}\label{q8:chi}
\frac{(n-1)S^2}{\sigma^2}\sim\chi^2_{n-1},\qquad \E(\chi^2_\nu)=\nu,\quad \Var(\chi^2_\nu)=2\nu .
\end{equation}

\passo{Moments of $S^2$}{Writing $S^2=\frac{\sigma^2}{n-1}\cdot\frac{(n-1)S^2}{\sigma^2}$ and using \eqref{q8:chi}:}
\begin{align}
\E[S^2]&=\frac{\sigma^2}{n-1}\,(n-1)=\sigma^2, \label{q8:es2}\\
\Var[S^2]&=\frac{\sigma^4}{(n-1)^2}\,\big(2(n-1)\big)=\frac{2\sigma^4}{n-1}. \label{q8:vs2}
\end{align}

\passo{(1.3.10)}{The decomposition}
\begin{equation}\label{q8:mse}
EQM[T]=\Var[T]+\big(\E[T]-\sigma^2\big)^2
\end{equation}
with \eqref{q8:es2} (zero bias) and \eqref{q8:vs2} substituted gives
$EQM[S^2]=\frac{2\sigma^4}{n-1}+(0)^2=\frac{2\sigma^4}{n-1}$, which is (1.3.10).

\passo{Moments of $\hat\sigma^2$}{Since}
\begin{equation}\label{q8:rel}
\hat\sigma^2=\frac{n-1}{n}S^2 ,
\end{equation}
substituting \eqref{q8:es2} and \eqref{q8:vs2}:
\begin{align}
\E[\hat\sigma^2]&=\frac{n-1}{n}\,(\sigma^2)\;\Longrightarrow\;\text{bias}=\frac{n-1}{n}\sigma^2-\sigma^2=-\frac{\sigma^2}{n}, \label{q8:bias}\\
\Var[\hat\sigma^2]&=\frac{(n-1)^2}{n^2}\Big(\frac{2\sigma^4}{n-1}\Big)=\frac{2(n-1)\sigma^4}{n^2}. \label{q8:vhat}
\end{align}

\passo{(1.3.11)}{Substituting \eqref{q8:bias} and \eqref{q8:vhat} into \eqref{q8:mse}:}
\begin{equation}\label{q8:msehat}
EQM[\hat\sigma^2]=\frac{2(n-1)\sigma^4}{n^2}+\Big(-\frac{\sigma^2}{n}\Big)^2=\frac{(2n-1)\sigma^4}{n^2}.
\end{equation}
Expanding the book's form:
\[
\frac{2\sigma^4}{n-1}\left[1-\frac{3n-1}{2n^2}\right]=\frac{2\sigma^4}{n-1}\cdot\frac{2n^2-3n+1}{2n^2}
=\frac{2\sigma^4}{n-1}\cdot\frac{(2n-1)(n-1)}{2n^2}=\frac{(2n-1)\sigma^4}{n^2},
\]
equal to \eqref{q8:msehat}, which verifies (1.3.11).

\passo{Comparison}{The bracket in (1.3.11) is below 1 for every $n\ge1$, so
$EQM[\hat\sigma^2]<EQM[S^2]$. For $n=10$: $EQM[S^2]=2/9\,\sigma^4=0.2222\,\sigma^4$ and
$EQM[\hat\sigma^2]=19/100\,\sigma^4=0.19\,\sigma^4$.}

\intuicao{Multiplying $S^2$ by $(n-1)/n<1$ adds a small downward bias but cuts the variance by more, so
the biased MLE wins on MSE. Unbiasedness is one criterion among several, not a guarantee of accuracy.}
\refbib{Bolfarine \& Sandoval (2010), §1.3, Exemplos 1.3.3 and 1.3.6 (pp.~6--9).}
\end{sol}

\questao{B 1.9 --- Bernoulli: unbiased $X$ versus the constant $1/2$.}
\fonte{Bolfarine, H. and Sandoval, M.~C. (2010), \textit{Introdução à Inferência Estatística},
§1.4 Exercícios, Exercício~1.9, p.~14 --- reproduced verbatim.}

\begin{otherlanguage}{brazil}
\textbf{1.9.} Seja $X$ uma única variável aleatória com distribuição de Bernoulli com parâmetro
$\theta$. Sejam $\hat\theta_1=X$ e $\hat\theta_2=1/2$ dois estimadores de $\theta$.\\
(i) Verifique se $\hat\theta_1$ e $\hat\theta_2$ são não viciados para $\theta$.\\
(ii) Compare os $EQMs$. Faça um gráfico dos $EQMs$ como função de $\theta$.
\end{otherlanguage}
\resolespaco
\begin{sol}
\solhead
\begin{notacao}
\item $X\sim\text{Bernoulli}(\theta)$, $\theta\in(0,1)$: $P(X=1)=\theta$, $\E X=\theta$, $\Var X=\theta(1-\theta)$.
\item $EQM_j(\theta)=\E[(\hat\theta_j-\theta)^2]$, $j=1,2$.
\end{notacao}

\passo{(i) Bias}{}
\begin{equation}\label{q9:bias}
\E[\hat\theta_1]=\E[X]=\theta,\qquad \E[\hat\theta_2]=\tfrac12 .
\end{equation}
$\hat\theta_1$ is unbiased; $\hat\theta_2$ has bias $\tfrac12-\theta$, so it is unbiased only at $\theta=1/2$.

\passo{(ii) MSEs}{Using $EQM=\Var+\text{bias}^2$ with \eqref{q9:bias}:}
\begin{align}
EQM_1(\theta)&=\theta(1-\theta)+0^2=\theta(1-\theta), \label{q9:m1}\\
EQM_2(\theta)&=0+\big(\tfrac12-\theta\big)^2=\big(\tfrac12-\theta\big)^2. \label{q9:m2}
\end{align}

\passo{(ii) Where each wins}{$EQM_2<EQM_1$ means, by \eqref{q9:m1}--\eqref{q9:m2},}
\[
\tfrac14-\theta+\theta^2<\theta-\theta^2 .
\]
Collecting terms:
\begin{equation}\label{q9:quad}
2\theta^2-2\theta+\tfrac14<0 .
\end{equation}
The roots of \eqref{q9:quad} are
\[
\theta=\frac{2\pm\sqrt{4-4\cdot2\cdot\frac14}}{2\cdot2}=\frac{2\pm\sqrt2}{4},
\quad\text{i.e.}\quad 0.1464\ \text{and}\ 0.8536 .
\]
So $\hat\theta_2=1/2$ has smaller MSE for $\theta\in(0.1464,\,0.8536)$ and $\hat\theta_1=X$ for
$\theta$ outside it; the MSEs are equal at the two roots.
\begin{center}
\begin{tikzpicture}
\begin{axis}[width=0.55\linewidth,height=5cm,xlabel={$\theta$},ylabel={$EQM$},xmin=0,xmax=1,
  ymin=0,ymax=0.27,samples=101,domain=0:1,legend pos=outer north east,legend style={font=\small},
  every axis plot/.append style={thick}]
\addplot[blue!60!black] {x*(1-x)}; \addlegendentry{$EQM_1=\theta(1-\theta)$}
\addplot[red!70!black,dashed] {(0.5-x)^2}; \addlegendentry{$EQM_2=(1/2-\theta)^2$}
\addplot[gray,densely dotted] coordinates {(0.1464,0) (0.1464,0.125)};
\addplot[gray,densely dotted] coordinates {(0.8536,0) (0.8536,0.125)};
\end{axis}
\end{tikzpicture}
\end{center}

\intuicao{A constant guess has zero variance and is excellent when it happens to be near the truth,
and poor otherwise. Neither estimator has smaller MSE for every $\theta$, so MSE alone cannot rank
them; this is why the theory restricts the class (unbiased estimators) before looking for a best one.}
\refbib{Bolfarine \& Sandoval (2010), §1.3 (erro quadrático médio, pp.~4--13).}
\end{sol}

\questao{B 1.12 --- $U(0,\theta)$: $c_1\overline X$ versus $c_2X_{(n)}$.}
\fonte{Bolfarine, H. and Sandoval, M.~C. (2010), \textit{Introdução à Inferência Estatística},
§1.4 Exercícios, Exercício~1.12, p.~14 --- reproduced verbatim.}

\begin{otherlanguage}{brazil}
\textbf{1.12.} Sejam $X_1,\dots,X_n$ uma amostra aleatória de tamanho $n$ da distribuição de uma
variável aleatória $X\sim U(0,\theta)$. Considere os estimadores $\hat\theta_1=c_1\overline X$ e
$\hat\theta_2=c_2X_{(n)}$.\\
(i) Encontre $c_1$ e $c_2$ que tornam os estimadores não viciados.\\
(ii) Encontre e compare os $EQMs$ dos dois estimadores.
\end{otherlanguage}
\resolespaco
\begin{sol}
\solhead
\begin{notacao}
\item $X_i$ i.i.d.\ $U(0,\theta)$, $\theta>0$: $\E X=\theta/2$, $\Var X=\theta^2/12$.
\item $X_{(n)}=\max_iX_i$; $c_1,c_2>0$ constants; $EQM_j=\E[(\hat\theta_j-\theta)^2]$.
\end{notacao}

\passo{Distribution of $X_{(n)}$}{$X_{(n)}\le x$ iff every $X_i\le x$, so for $0\le x\le\theta$}
\begin{equation}\label{q10:cdf}
F_{(n)}(x)=\Big(\frac x\theta\Big)^n,\qquad f_{(n)}(x)=\frac{n\,x^{n-1}}{\theta^n}
\end{equation}
(Bolfarine's (1.3.14)). Its moments follow from \eqref{q10:cdf}:
\begin{equation}\label{q10:mom}
\E[X_{(n)}^r]=\int_0^\theta x^r\frac{nx^{n-1}}{\theta^n}dx=\frac{n\,\theta^r}{n+r}.
\end{equation}
Setting $r=1,2$ in \eqref{q10:mom}:
\begin{equation}\label{q10:m12}
\E[X_{(n)}]=\frac{n\theta}{n+1},\qquad \E[X_{(n)}^2]=\frac{n\theta^2}{n+2}.
\end{equation}
Hence
\begin{equation}\label{q10:var}
\Var[X_{(n)}]=\frac{n\theta^2}{n+2}-\Big(\frac{n\theta}{n+1}\Big)^2
=\frac{n\theta^2\big[(n+1)^2-n(n+2)\big]}{(n+1)^2(n+2)}=\frac{n\theta^2}{(n+1)^2(n+2)}
\end{equation}
(Bolfarine's (1.3.15)).

\passo{(i) Unbiasing constants}{$\E[c_1\overline X]=c_1\theta/2=\theta$ and, by \eqref{q10:m12},
$\E[c_2X_{(n)}]=c_2\,n\theta/(n+1)=\theta$:}
\begin{equation}\label{q10:c}
c_1=2,\qquad c_2=\frac{n+1}{n}.
\end{equation}

\passo{(ii) MSE of $\hat\theta_1$}{Both are unbiased, so $EQM=\Var$. With $\Var\overline X=\theta^2/(12n)$:}
\begin{equation}\label{q10:m1}
EQM_1=\Var(2\overline X)=4\cdot\frac{\theta^2}{12n}=\frac{\theta^2}{3n}.
\end{equation}

\passo{(ii) MSE of $\hat\theta_2$}{Substituting \eqref{q10:c} and \eqref{q10:var}:}
\begin{equation}\label{q10:m2}
EQM_2=\Big(\frac{n+1}{n}\Big)^2\cdot\frac{n\theta^2}{(n+1)^2(n+2)}=\frac{\theta^2}{n(n+2)}.
\end{equation}

\passo{(ii) Comparison}{Dividing \eqref{q10:m2} by \eqref{q10:m1}:}
\begin{equation}\label{q10:ratio}
\frac{EQM_2}{EQM_1}=\frac{\theta^2/[n(n+2)]}{\theta^2/(3n)}=\frac{3}{n+2},
\end{equation}
which is $1$ at $n=1$ (both estimators are $2X_1$) and below 1 for every $n\ge2$. For $n=10$,
$\theta=1$: $EQM_1=1/30=0.0333$, $EQM_2=1/120=0.0083$, ratio $3/12=0.25$. A Monte Carlo run with
200{,}000 samples (check script) gives $0.0333$ and $0.0084$.

\intuicao{$X_{(n)}$ is sufficient for $\theta$ and sits right below the upper end of the support,
which is exactly what we estimate; $\overline X$ throws that information away. Because the support
depends on $\theta$, the MSE of $\hat\theta_2$ falls like $1/n^2$, faster than the $1/n$ of regular
problems (Cramér--Rao does not apply).}
\refbib{Bolfarine \& Sandoval (2010), §1.3, Exemplo 1.3.8 and (1.3.14)--(1.3.15) (pp.~10--11).}
\end{sol}

\questao{B 2.2 --- Binomial$(2,\theta)$: Cramér--Rao bound and an efficient estimator.}
\fonte{Bolfarine, H. and Sandoval, M.~C. (2010), \textit{Introdução à Inferência Estatística},
§2.6 Exercícios, Exercício~2.2, p.~32 --- reproduced verbatim.}

\begin{otherlanguage}{brazil}
\textbf{2.2.} Sejam $X_1,\dots,X_n$ uma amostra aleatória da variável aleatória
$X\sim\mathit{Binomial}(2,\theta)$.\\
(i) Encontre o limite inferior da variância dos estimadores não viciados de $\theta$.\\
(ii) Encontre uma estatística suficiente para $\theta$.\\
(iii) Obtenha um estimador não viciado para $\theta$ que seja função da estatística suficiente.\\
(iv) Verifique se o estimador é eficiente.
\end{otherlanguage}
\resolespaco
\begin{sol}
\solhead
\begin{notacao}
\item $f(x\mid\theta)=\binom2x\theta^x(1-\theta)^{2-x}$, $x\in\{0,1,2\}$, $\theta\in(0,1)$;
  $\E X=2\theta$, $\Var X=2\theta(1-\theta)$.
\item $LI(\theta)$: Cramér--Rao lower bound, Bolfarine (2.1.1); $T=\sum_iX_i\sim\text{Binomial}(2n,\theta)$.
\end{notacao}

\passo{(i) Score}{}
\begin{equation}\label{q11:logf}
\log f(x\mid\theta)=\log\tbinom2x+x\log\theta+(2-x)\log(1-\theta).
\end{equation}
Differentiating \eqref{q11:logf}:
\[
\frac{\partial\log f}{\partial\theta}=\frac x\theta-\frac{2-x}{1-\theta}
=\frac{x(1-\theta)-(2-x)\theta}{\theta(1-\theta)} .
\]
Simplifying the numerator:
\begin{equation}\label{q11:score}
\frac{\partial\log f}{\partial\theta}=\frac{x-2\theta}{\theta(1-\theta)}.
\end{equation}

\passo{(i) Information and bound}{Squaring \eqref{q11:score} and taking expectations, with $\E(X-2\theta)^2=\Var X$:}
\begin{equation}\label{q11:info}
\E\Big[\Big(\frac{\partial\log f(X\mid\theta)}{\partial\theta}\Big)^2\Big]=\frac{2\theta(1-\theta)}{\theta^2(1-\theta)^2}=\frac{2}{\theta(1-\theta)}.
\end{equation}
Substituting \eqref{q11:info} into Bolfarine's (2.1.1), $LI(\theta)=1/\{nE[(\partial\log f/\partial\theta)^2]\}$:
\begin{equation}\label{q11:LI}
LI(\theta)=\frac{1}{n\cdot\frac{2}{\theta(1-\theta)}}=\frac{\theta(1-\theta)}{2n}.
\end{equation}

\passo{(ii) Sufficient statistic}{The joint probability function is}
\begin{equation}\label{q11:joint}
\prod_{i=1}^nf(x_i\mid\theta)=\underbrace{\prod_{i=1}^n\tbinom2{x_i}}_{h(\mathbf x)}\;
\underbrace{\theta^{\sum x_i}(1-\theta)^{2n-\sum x_i}}_{g_\theta(T(\mathbf x))},
\end{equation}
so by the factorisation criterion (Teorema 2.2.1) $T=\sum_{i=1}^nX_i$ is sufficient.

\passo{(iii) Unbiased function of $T$}{$\E T=n\,\E X=2n\theta$, so}
\begin{equation}\label{q11:est}
\hat\theta=\frac{T}{2n}=\frac{\overline X}{2},\qquad \E\hat\theta=\frac{2n\theta}{2n}=\theta .
\end{equation}

\passo{(iv) Efficiency}{With $\Var T=2n\theta(1-\theta)$:}
\begin{equation}\label{q11:var}
\Var\hat\theta=\frac{\Var T}{4n^2}=\frac{2n\theta(1-\theta)}{4n^2}=\frac{\theta(1-\theta)}{2n}.
\end{equation}
\eqref{q11:var} equals \eqref{q11:LI}, so $e(\hat\theta)=LI(\theta)/\Var\hat\theta=1$: $\hat\theta$ is
efficient (and, $T$ being complete, also the ENVVUM).

\intuicao{A Binomial$(2,\theta)$ observation is two Bernoulli trials, so the sample is $2n$ trials and
the natural estimator is the success proportion $T/(2n)$. The score \eqref{q11:score} is linear in
$x$, which is the condition for the bound to be attained.}
\refbib{Bolfarine \& Sandoval (2010), §2.1 (Teorema 2.1.1, pp.~15--19) and §2.2 (Teorema 2.2.1, pp.~20--23).}
\end{sol}

\questao{B 2.4 --- $X_1+2X_2$ is not sufficient for a Poisson mean.}
\fonte{Bolfarine, H. and Sandoval, M.~C. (2010), \textit{Introdução à Inferência Estatística},
§2.6 Exercícios, Exercício~2.4, p.~32 --- reproduced verbatim.}

\begin{otherlanguage}{brazil}
\textbf{2.4.} Sejam $X_1,X_2$ uma amostra aleatória da variável aleatória $X\sim\mathit{Poisson}(\theta)$.
Mostre que $T=X_1+2X_2$ não é suficiente para $\theta$.
\end{otherlanguage}
\resolespaco
\begin{sol}
\solhead
\begin{notacao}
\item $X_1,X_2$ i.i.d.\ Poisson$(\theta)$, $P(X=x)=\theta^xe^{-\theta}/x!$, $\theta>0$; $T=X_1+2X_2$.
\item Sufficiency (Definição 2.2.1): $P(\mathbf X=\mathbf x\mid T=t)$ does not depend on $\theta$.
\end{notacao}

\passo{Pick a value of $T$}{The pairs with $T=2$ are}
\begin{equation}\label{q12:set}
\{T=2\}=\{(X_1,X_2)=(2,0)\}\cup\{(X_1,X_2)=(0,1)\}.
\end{equation}

\passo{Probabilities}{By independence,}
\begin{align}
P(X_1=2,X_2=0)&=\frac{\theta^2e^{-\theta}}{2!}\cdot e^{-\theta}=\frac{\theta^2e^{-2\theta}}{2}, \label{q12:p20}\\
P(X_1=0,X_2=1)&=e^{-\theta}\cdot\theta e^{-\theta}=\theta e^{-2\theta}. \label{q12:p01}
\end{align}
Adding \eqref{q12:p20} and \eqref{q12:p01} as \eqref{q12:set} requires:
\begin{equation}\label{q12:pT}
P(T=2)=e^{-2\theta}\Big(\frac{\theta^2}{2}+\theta\Big).
\end{equation}

\passo{Conditional probability}{Dividing \eqref{q12:p20} by \eqref{q12:pT}:}
\[
P(X_1=2,X_2=0\mid T=2)=\frac{\theta^2e^{-2\theta}/2}{e^{-2\theta}(\theta^2/2+\theta)}=\frac{\theta^2/2}{\theta^2/2+\theta}.
\]
Dividing numerator and denominator by $\theta/2$:
\begin{equation}\label{q12:cond}
P(X_1=2,X_2=0\mid T=2)=\frac{\theta}{\theta+2},
\end{equation}
which depends on $\theta$ (e.g.\ $1/3=0.3333$ at $\theta=1$, $2/4=0.5$ at $\theta=2$). So $T$ is not
sufficient.

\intuicao{The sufficient statistic here is $X_1+X_2$. Weighting $X_2$ twice merges sample points with
different totals: $T=2$ mixes ``two events'' with ``one event'', and the split between these cases
still carries information about $\theta$.}
\refbib{Bolfarine \& Sandoval (2010), §2.2 (estatísticas suficientes, pp.~20--23; Exemplo 2.2.4).}
\end{sol}

\questao{B 2.10 --- ENVVUM of $\theta(1-\theta)$ for Bernoulli data.}
\fonte{Bolfarine, H. and Sandoval, M.~C. (2010), \textit{Introdução à Inferência Estatística},
§2.6 Exercícios, Exercício~2.10, p.~33 --- reproduced verbatim.}

\begin{otherlanguage}{brazil}
\textbf{2.10.} Sejam $X_1,\dots,X_n$ uma amostra aleatória da variável aleatória.
$X\sim\mathit{Bernoulli}(\theta)$. Obtenha o ENVVUM para $\theta(1-\theta)$.\\
Sugestão: verifique se $S^2=\frac{n}{(n-1)}\overline X(1-\overline X)$ é não viciado para $\theta(1-\theta)$.
\end{otherlanguage}
\resolespaco
\begin{sol}
\solhead
\begin{notacao}
\item $X_i$ i.i.d.\ Bernoulli$(\theta)$, $\theta\in(0,1)$, $n\ge2$; $T=\sum_iX_i\sim\text{Binomial}(n,\theta)$;
  $\overline X=T/n$.
\item ENVVUM: \emph{estimador não viciado de variância uniformemente mínima} (UMVUE).
\end{notacao}

\passo{Moments of $\overline X$}{}
\begin{equation}\label{q13:m}
\E\overline X=\theta,\qquad \Var\overline X=\frac{\theta(1-\theta)}{n}.
\end{equation}
Hence
\begin{equation}\label{q13:m2}
\E\overline X^2=\Var\overline X+(\E\overline X)^2=\frac{\theta(1-\theta)}{n}+\theta^2 .
\end{equation}

\passo{Bias of the plug-in}{Using \eqref{q13:m} and \eqref{q13:m2}:}
\[
\E\big[\overline X(1-\overline X)\big]=\E\overline X-\E\overline X^2=\theta-\Big(\frac{\theta(1-\theta)}{n}+\theta^2\Big).
\]
Since $\theta-\theta^2=\theta(1-\theta)$:
\begin{equation}\label{q13:bias}
\E\big[\overline X(1-\overline X)\big]=\theta(1-\theta)-\frac{\theta(1-\theta)}{n}=\frac{n-1}{n}\,\theta(1-\theta).
\end{equation}

\passo{Unbiased correction}{Multiplying \eqref{q13:bias} by $n/(n-1)$:}
\begin{equation}\label{q13:unb}
\E[S^2]=\frac{n}{n-1}\cdot\frac{n-1}{n}\,\theta(1-\theta)=\theta(1-\theta).
\end{equation}

\passo{Function of a complete sufficient statistic}{Substituting $\overline X=T/n$:}
\begin{equation}\label{q13:T}
S^2=\frac{n}{n-1}\cdot\frac Tn\Big(1-\frac Tn\Big)=\frac{T(n-T)}{n(n-1)} .
\end{equation}
$T$ is sufficient and complete for the Bernoulli family (one-parameter exponential family, Teorema
2.5.2). By Lehmann--Scheffé (Teorema 2.5.3), the unbiased estimator \eqref{q13:unb} that is a function
of $T$, \eqref{q13:T}, is the unique ENVVUM of $\theta(1-\theta)$.

\intuicao{Since $X_i^2=X_i$, $S^2$ is exactly the usual sample variance with divisor $n-1$. The plug-in
$\overline X(1-\overline X)$ is the MLE of $\theta(1-\theta)$ (Exemplo 3.2.1) and is biased downward by
the factor $(n-1)/n$, the same correction as $\hat\sigma^2$ versus $S^2$ in B~1.4.}
\refbib{Bolfarine \& Sandoval (2010), §2.5 (completude, Teoremas 2.5.2--2.5.3, pp.~29--32).}
\end{sol}

\questao{B 3.6 --- MLE of $\mu^2$ versus the unbiased $\overline X^2-1/n$.}
\fonte{Bolfarine, H. and Sandoval, M.~C. (2010), \textit{Introdução à Inferência Estatística},
§3.8 Exercícios, Exercício~3.6, p.~53 --- reproduced verbatim. Exercise 2.9, to which it refers, is
reproduced verbatim from §2.6, p.~33.}

\begin{otherlanguage}{brazil}
\textbf{3.6.} Encontre o estimador de máxima verossimilhança de $\theta^2$ no Exercício 2.9 e compare
seu erro quadrático médio com o do estimador eficiente $\hat\gamma$ dado no Exercício 2.9, (i).

\smallskip
\emph{Exercício referido (p.~33):} \textbf{2.9.} Sejam $X_1,\dots,X_n$ uma amostra aleatória da
variável aleatória $X\sim N(\mu,1)$.\\
(i) Mostre que o estimador $\hat\gamma=\overline X^2-1/n$ é não viciado para $g(\mu)=\mu^2$.\\
(ii) Existe ENVVUM para $\mu^2$?\\
(iii) Encontre o limite inferior da variância dos estimadores não viciados de $g(\mu)=\mu^2$ e
verifique se $\hat\gamma$ é eficiente.
\end{otherlanguage}

\noindent{\small\emph{Note:} the parameter of Exercise 2.9 is $\mu$; the ``$\theta^2$'' of 3.6 is read as
$\mu^2$.}
\resolespaco
\begin{sol}
\solhead
\begin{notacao}
\item $X_i$ i.i.d.\ $N(\mu,1)$, $\mu\in\mathbb R$; $\overline X\sim N(\mu,1/n)$; $Z\sim N(0,1)$.
\item $\widehat{\mu^2}$: MLE of $\mu^2$; $\hat\gamma=\overline X^2-1/n$; $EQM(T)=\E[(T-\mu^2)^2]$.
\end{notacao}

\passo{MLE}{The log-likelihood $\ell(\mu)=-\frac n2\log2\pi-\frac12\sum_i(x_i-\mu)^2$ has score
$\sum_i(x_i-\mu)=0$, so $\hat\mu=\overline X$. The map $\mu\mapsto\mu^2$ is not one-to-one, but the
likelihood of $\eta=\mu^2$, $L^*(\eta)=\max\{L(\sqrt\eta),L(-\sqrt\eta)\}$, is maximised at
$\eta=\overline X^2$ (the extension of the invariance principle, Teorema 3.2.2):}
\begin{equation}\label{q14:mle}
\widehat{\mu^2}=\overline X^2 .
\end{equation}

\passo{Moments of $\overline X^2$}{Write}
\begin{equation}\label{q14:Z}
\overline X=\mu+\frac{Z}{\sqrt n}\;\Longrightarrow\;\overline X^2=\mu^2+\frac{2\mu Z}{\sqrt n}+\frac{Z^2}{n}.
\end{equation}
With $\E Z=0$, $\E Z^2=1$, $\Var Z^2=2$, $\Cov(Z,Z^2)=\E Z^3=0$, \eqref{q14:Z} gives
\begin{align}
\E\overline X^2&=\mu^2+\frac1n, \label{q14:E}\\
\Var\overline X^2&=\frac{4\mu^2}{n}\Var Z+\frac{1}{n^2}\Var Z^2=\frac{4\mu^2}{n}+\frac{2}{n^2}. \label{q14:V}
\end{align}

\passo{MSE of the MLE}{By \eqref{q14:E} the bias is $1/n$; adding its square to \eqref{q14:V}:}
\begin{equation}\label{q14:mse1}
EQM(\overline X^2)=\frac{4\mu^2}{n}+\frac{2}{n^2}+\Big(\frac1n\Big)^2=\frac{4\mu^2}{n}+\frac{3}{n^2}.
\end{equation}

\passo{MSE of $\hat\gamma$}{By \eqref{q14:E}, $\E\hat\gamma=\mu^2$; subtracting a constant leaves the
variance \eqref{q14:V} unchanged:}
\begin{equation}\label{q14:mse2}
EQM(\hat\gamma)=\Var\hat\gamma=\frac{4\mu^2}{n}+\frac{2}{n^2}.
\end{equation}

\passo{Comparison}{Subtracting \eqref{q14:mse2} from \eqref{q14:mse1}:}
\[
EQM(\overline X^2)-EQM(\hat\gamma)=\frac{3}{n^2}-\frac{2}{n^2}=\frac1{n^2}>0 ,
\]
so $\hat\gamma$ has smaller MSE for every $\mu$ and $n$. For $\mu=1$, $n=5$: $0.92$ versus $0.88$
(Monte Carlo with 400{,}000 replicates reproduces both).

\passo{A remark on ``eficiente''}{$\hat\gamma$ is a function of the complete sufficient statistic
$\overline X$ and unbiased, so it is the ENVVUM (the answer to 2.9(ii)). It is \emph{not} efficient in
the Cramér--Rao sense: the bound for $g(\mu)=\mu^2$ is $(g'(\mu))^2/n=4\mu^2/n$, strictly below
\eqref{q14:mse2} (answer to 2.9(iii)). The word ``eficiente'' in 3.6 should be read as ``best unbiased''.}

\intuicao{Squaring a noisy $\overline X$ adds the noise variance $1/n$ on average, so $\overline X^2$
overestimates $\mu^2$; $\hat\gamma$ removes exactly that bias without adding variance. Invariance
gives the MLE for free but says nothing about its finite-sample bias. ($\hat\gamma$ can be negative;
$\max\{\hat\gamma,0\}$ would lower the MSE further.)}
\refbib{Bolfarine \& Sandoval (2010), §3.2.1 (invariância, Teorema 3.2.2, p.~43) and §2.1 (pp.~15--19).}
\end{sol}

\questao{B 3.13 --- MLE of $P[X>1]$ for exponential data and its large-sample distribution.}
\fonte{Bolfarine, H. and Sandoval, M.~C. (2010), \textit{Introdução à Inferência Estatística},
§3.8 Exercícios, Exercício~3.13, p.~54 --- reproduced verbatim.}

\begin{otherlanguage}{brazil}
\textbf{3.13.} Sejam $X_1,\dots,X_n$ uma amostra aleatória da variável aleatória $X$ com distribuição
exponencial com parâmetro $\theta$. Encontre o estimador de máxima verossimilhança de $g(\theta)=P[X>1]$
e sua distribuição aproximada quando $n$ for grande.
\end{otherlanguage}
\resolespaco
\begin{sol}
\solhead
\begin{notacao}
\item $f(x\mid\theta)=\theta e^{-\theta x}$, $x>0$, $\theta>0$ (Bolfarine's parametrisation, Exemplo
  3.2.3): $\E X=1/\theta$.
\item $g(\theta)=P[X>1]=e^{-\theta}$; $I_F(\theta)$: Fisher information per observation.
\end{notacao}

\passo{MLE of $\theta$}{}
\begin{equation}\label{q15:ll}
\ell(\theta)=n\log\theta-\theta\sum_{i=1}^nx_i .
\end{equation}
Differentiating \eqref{q15:ll} and setting the score to zero:
\begin{equation}\label{q15:mle}
\ell'(\theta)=\frac n\theta-\sum_ix_i=0\;\Longrightarrow\;\hat\theta=\frac{1}{\overline X},
\end{equation}
a maximum since $\ell''(\theta)=-n/\theta^2<0$.

\passo{MLE of $g$}{$g(\theta)=e^{-\theta}$ is one-to-one on $\theta>0$, so by invariance (Teorema 3.2.2)
and \eqref{q15:mle}:}
\begin{equation}\label{q15:g}
g(\hat\theta)=e^{-1/\overline X}.
\end{equation}

\passo{Fisher information}{$\log f=\log\theta-\theta x$, so $\partial^2\log f/\partial\theta^2=-1/\theta^2$ and}
\begin{equation}\label{q15:I}
I_F(\theta)=-\E\Big[\frac{\partial^2\log f(X\mid\theta)}{\partial\theta^2}\Big]=\frac1{\theta^2}.
\end{equation}

\passo{Delta method}{$g'(\theta)=-e^{-\theta}$. Substituting it and \eqref{q15:I} into Bolfarine's (3.2.4),
$\sqrt n\,(g(\hat\theta)-g(\theta))\overset{a}{\sim}N\big(0,(g'(\theta))^2/I_F(\theta)\big)$:}
\[
\sqrt n\,\big(e^{-1/\overline X}-e^{-\theta}\big)\overset{a}{\sim}N\Big(0,\ \frac{(-e^{-\theta})^2}{1/\theta^2}\Big).
\]
Simplifying the variance:
\begin{equation}\label{q15:asy}
e^{-1/\overline X}\overset{a}{\sim}N\Big(e^{-\theta},\ \frac{\theta^2e^{-2\theta}}{n}\Big).
\end{equation}

\intuicao{The MLE of a probability is the probability evaluated at the MLE, and the delta method turns
the $1/n$ variance of $\hat\theta$ into a variance for $g(\hat\theta)$ scaled by $g'(\theta)^2$. The
model-free alternative, the share of $X_i>1$, has variance $e^{-\theta}(1-e^{-\theta})/n$, never smaller
than \eqref{q15:asy}: using the exponential model buys precision if the model is right.}
\refbib{Bolfarine \& Sandoval (2010), §3.2.1--3.2.2, Exemplo 3.2.3 and (3.2.4) (p.~44).}
\end{sol}

\questao{B 3.17 --- consistency of the estimators in five examples.}
\fonte{Bolfarine, H. and Sandoval, M.~C. (2010), \textit{Introdução à Inferência Estatística},
§3.8 Exercícios, Exercício~3.17, p.~55 --- reproduced verbatim. The list of the examples' estimators
below the statement is a summary (not verbatim) of pp.~37, 43, 44 and 51.}

\begin{otherlanguage}{brazil}
\textbf{3.17.} Verifique se os estimadores obtidos nos Exemplos 3.1.2, 3.1.3, 3.2.1, 3.2.3 e 3.6.2
são consistentes.
\end{otherlanguage}

\smallskip
\noindent{\small\emph{Estimators in the examples (summary):} 3.1.2, Bernoulli$(\theta)$:
$\hat\theta=\overline X$; 3.1.3, $U(0,\theta)$: $\hat\theta=X_{(n)}$; 3.2.1, Bernoulli$(\theta)$:
$g(\hat\theta)=\overline X(1-\overline X)$ for $\theta(1-\theta)$; 3.2.3, $Exp(\theta)$ with density
$\theta e^{-\theta x}$: $g(\hat\theta)=e^{-1/\overline X}$ for $P_\theta[X>1]=e^{-\theta}$; 3.6.2,
gamma$(\alpha,\beta)$ with density $\beta^\alpha x^{\alpha-1}e^{-\beta x}/\Gamma(\alpha)$, method of
moments: $\hat\alpha=\overline X^2/\hat\sigma^2$, $\hat\beta=\overline X/\hat\sigma^2$,
$\hat\sigma^2=\sum_i(X_i-\overline X)^2/n$.}
\resolespaco
\begin{sol}
\solhead
\begin{notacao}
\item Consistency (Definição 3.7.1): $\lim_{n\to\infty}P(|\hat\theta-\theta|>\epsilon)=0$ for every
  $\epsilon>0$, written $\hat\theta\xrightarrow{p}\theta$.
\item $m_2=\frac1n\sum_iX_i^2$; $EQM(\hat\theta)=\E[(\hat\theta-\theta)^2]$.
\end{notacao}

\passo{Tools}{Markov's inequality applied to $(\hat\theta-\theta)^2$ (Chebyshev, as in Exemplo 3.7.1):}
\begin{equation}\label{q16:cheb}
P(|\hat\theta-\theta|>\epsilon)\le\frac{EQM(\hat\theta)}{\epsilon^2},
\end{equation}
so $EQM(\hat\theta)\to0$ implies consistency. And the continuous-mapping theorem (standard; not
stated in Bolfarine): if $\hat\theta\xrightarrow{p}\theta$ and $h$ is continuous at $\theta$, then
\begin{equation}\label{q16:cmt}
h(\hat\theta)\xrightarrow{p}h(\theta),
\end{equation}
also for vectors $\hat\theta$.

\passo{Exemplo 3.1.2: $\overline X$}{Unbiased with $EQM=\Var\overline X=\theta(1-\theta)/n\to0$, so by
\eqref{q16:cheb} $\overline X\xrightarrow{p}\theta$. Consistent.}

\passo{Exemplo 3.1.3: $X_{(n)}$}{Since $X_{(n)}\le\theta$, for $0<\epsilon<\theta$ and by
\eqref{q10:cdf}:}
\begin{equation}\label{q16:max}
P(|X_{(n)}-\theta|>\epsilon)=P(X_{(n)}<\theta-\epsilon)=\Big(\frac{\theta-\epsilon}{\theta}\Big)^n\xrightarrow[n\to\infty]{}0 ,
\end{equation}
and for $\epsilon\ge\theta$ the probability is 0. Consistent (although biased, Exemplo 1.3.8).

\passo{Exemplo 3.2.1: $\overline X(1-\overline X)$}{$\overline X\xrightarrow{p}\theta$ (Exemplo 3.1.2) and
$h(t)=t(1-t)$ is continuous, so by \eqref{q16:cmt} $\overline X(1-\overline X)\xrightarrow{p}\theta(1-\theta)$.
Consistent; its bias $-\theta(1-\theta)/n$ from \eqref{q13:bias} vanishes.}

\passo{Exemplo 3.2.3: $e^{-1/\overline X}$}{$\E X=1/\theta$ and $\Var X=1/\theta^2$, so
$EQM(\overline X)=1/(n\theta^2)\to0$ and, by \eqref{q16:cheb},}
\begin{equation}\label{q16:exp}
\overline X\xrightarrow{p}\frac1\theta .
\end{equation}
$h(t)=e^{-1/t}$ is continuous at $t=1/\theta>0$; by \eqref{q16:cmt} and \eqref{q16:exp},
$e^{-1/\overline X}\xrightarrow{p}e^{-\theta}$. Consistent.

\passo{Exemplo 3.6.2: $\hat\alpha,\hat\beta$}{For the gamma, $\E X=\alpha/\beta$,
$\E X^2=\alpha/\beta^2+\alpha^2/\beta^2$, and all moments are finite, so by \eqref{q16:cheb}
(the variances of $\overline X$ and $m_2$ are $O(1/n)$):}
\begin{equation}\label{q16:mom}
\overline X\xrightarrow{p}\frac\alpha\beta,\qquad m_2\xrightarrow{p}\frac{\alpha}{\beta^2}+\frac{\alpha^2}{\beta^2}.
\end{equation}
Since $\hat\sigma^2=m_2-\overline X^2$, \eqref{q16:cmt} applied to \eqref{q16:mom} gives
\begin{equation}\label{q16:s2}
\hat\sigma^2\xrightarrow{p}\Big(\frac{\alpha}{\beta^2}+\frac{\alpha^2}{\beta^2}\Big)-\Big(\frac\alpha\beta\Big)^2=\frac{\alpha}{\beta^2}.
\end{equation}
The maps $(a,b)\mapsto a^2/b$ and $(a,b)\mapsto a/b$ are continuous for $b>0$; by \eqref{q16:cmt} with
\eqref{q16:mom} and \eqref{q16:s2}:
\[
\hat\alpha=\frac{\overline X^2}{\hat\sigma^2}\xrightarrow{p}\frac{(\alpha/\beta)^2}{\alpha/\beta^2}=\alpha,
\qquad
\hat\beta=\frac{\overline X}{\hat\sigma^2}\xrightarrow{p}\frac{\alpha/\beta}{\alpha/\beta^2}=\beta .
\]
Consistent. All five estimators are consistent.

\intuicao{Consistency is a weak, large-sample property: biased estimators ($X_{(n)}$,
$\overline X(1-\overline X)$, $\hat\alpha$) can be consistent, and smooth functions of consistent
estimators inherit it. Sample moments converge by the law of large numbers, and everything built
continuously from them follows.}
\refbib{Bolfarine \& Sandoval (2010), §3.7 (estimadores consistentes, Definição 3.7.1 and Exemplo
3.7.1, p.~52); examples on pp.~37, 43, 44, 51.}
\end{sol}

% =============================================================================
\bloco{6}{Morettin \& Singer, \textit{Estatística e Ciência de Dados}}{MS 6.15}

\questao{MS 6.15 --- decomposition of the total sum of squares.}
\fonte{Morettin, P.~A. and Singer, J.~M. (2021), \textit{Estatística e Ciência de Dados} (versão
preliminar, abril 2021), §6.7 Exercícios, Exercício~15, p.~240 --- reproduced verbatim. The
definitions it uses are reproduced verbatim from §6.2, p.~185.}

\begin{otherlanguage}{brazil}
\textbf{15)} Mostre que $SQTot=SQRes+SQReg$.

\smallskip
\emph{Definições (§6.2, p.~185):} ``[\dots] em que a \textbf{soma de quadrados total} é
$SQTot=\sum_{i=1}^n(y_i-\overline y)^2$, a \textbf{soma de quadrados dos resíduos} é
$SQRes=\sum_{i=1}^n(y_i-\widehat y_i)^2$ e a \textbf{soma de quadrados da regressão} é
$SQReg=\sum_{i=1}^n(\widehat y_i-\overline y)^2$.'' (Model: $y_i=\alpha+\beta x_i+e_i$, with
$\widehat y_i=\widehat\alpha+\widehat\beta x_i$ fitted by least squares.)
\end{otherlanguage}
\resolespaco
\begin{sol}
\solhead
\begin{notacao}
\item $(x_i,y_i)$, $i=1,\dots,n$; $\widehat\alpha,\widehat\beta$: least-squares estimates (MS (6.6)--(6.7));
  $\widehat y_i=\widehat\alpha+\widehat\beta x_i$: fitted values; $\widehat e_i=y_i-\widehat y_i$: residuals;
  $\overline y$: sample mean of $y$.
\item $Q(\alpha,\beta)=\sum_i(y_i-\alpha-\beta x_i)^2$: least-squares criterion.
\end{notacao}

\passo{Normal equations}{Setting the derivatives of $Q$ to zero at $(\widehat\alpha,\widehat\beta)$:}
\begin{align}
\frac{\partial Q}{\partial\alpha}=-2\sum_i(y_i-\widehat\alpha-\widehat\beta x_i)=0&\;\Longrightarrow\;\sum_i\widehat e_i=0, \label{q17:ne1}\\
\frac{\partial Q}{\partial\beta}=-2\sum_ix_i(y_i-\widehat\alpha-\widehat\beta x_i)=0&\;\Longrightarrow\;\sum_ix_i\widehat e_i=0. \label{q17:ne2}
\end{align}

\passo{Decompose each deviation}{Adding and subtracting $\widehat y_i$:}
\begin{equation}\label{q17:dev}
y_i-\overline y=\widehat e_i+(\widehat y_i-\overline y).
\end{equation}
Squaring \eqref{q17:dev} and summing over $i$:
\begin{equation}\label{q17:sq}
\sum_i(y_i-\overline y)^2=\sum_i\widehat e_i^{\,2}+\sum_i(\widehat y_i-\overline y)^2+2\sum_i\widehat e_i(\widehat y_i-\overline y).
\end{equation}

\passo{The cross term vanishes}{Substituting $\widehat y_i=\widehat\alpha+\widehat\beta x_i$ in the cross term:}
\[
\sum_i\widehat e_i(\widehat y_i-\overline y)=\sum_i\widehat e_i\big((\widehat\alpha+\widehat\beta x_i)-\overline y\big)
=(\widehat\alpha-\overline y)\sum_i\widehat e_i+\widehat\beta\sum_ix_i\widehat e_i .
\]
Substituting \eqref{q17:ne1} and \eqref{q17:ne2}:
\begin{equation}\label{q17:cross}
\sum_i\widehat e_i(\widehat y_i-\overline y)=(\widehat\alpha-\overline y)\,(0)+\widehat\beta\,(0)=0 .
\end{equation}

\passo{Conclusion}{Substituting \eqref{q17:cross} into \eqref{q17:sq}:}
\[
\underbrace{\sum_i(y_i-\overline y)^2}_{SQTot}=\underbrace{\sum_i\widehat e_i^{\,2}}_{SQRes}+\underbrace{\sum_i(\widehat y_i-\overline y)^2}_{SQReg}.
\]
The same proof covers multiple regression with an intercept: the normal equations
$\mathbf X^\top\widehat{\mathbf e}=\mathbf 0$ include the column of ones (giving \eqref{q17:ne1}) and
imply $\widehat{\mathbf y}^\top\widehat{\mathbf e}=\widehat{\boldsymbol\beta}^\top\mathbf X^\top\widehat{\mathbf e}=0$.
Without an intercept \eqref{q17:ne1} fails in general and so does the identity. (The check script
verifies it numerically on simulated data.)

\intuicao{Geometrically, the residual vector is orthogonal to every column of the design, hence to
$\widehat{\mathbf y}-\overline y\mathbf 1$; the decomposition is Pythagoras' theorem in $\mathbb R^n$.
It is what makes $R^2=SQReg/SQTot$ a share between 0 and 1.}
\refbib{Morettin \& Singer (2021), §6.2 (pp.~184--185) and Figura 6.34 with the result on p.~227.}
\end{sol}

\end{document}
